\documentclass[11pt,a4paper,fontset=fandol]{ctexart}
\usepackage[left=19mm,right=19mm,top=18mm,bottom=18mm]{geometry}
\usepackage{amsmath,amssymb,bm,mathtools}
\usepackage{xcolor,graphicx,array,booktabs,enumitem,tikz}
\usetikzlibrary{arrows.meta,calc}
\definecolor{defblue}{HTML}{125889}
\definecolor{compgreen}{HTML}{23623B}
\definecolor{calcorange}{HTML}{974800}
\definecolor{warnred}{HTML}{A32830}
\definecolor{ink}{HTML}{18202B}
\color{ink}
\setlength{\parindent}{0pt}
\newcommand{\R}{\mathbb{R}}
\newcommand{\T}{^{\mathsf T}}
\newcommand{\norm}[1]{\lVert #1\rVert}
\newcommand{\sig}{\sigma}
\newcommand{\kw}[1]{\textbf{#1}}
\newcommand{\ans}[1]{\textbf{答卷答案：}#1}

\usepackage{fancyhdr,lastpage,needspace}
\usepackage[unicode,colorlinks=true,linkcolor=defblue,urlcolor=defblue,bookmarksnumbered=true]{hyperref}
\hypersetup{pdftitle={CS5486 期中 Cheatsheet 配套复习题与解析},pdfauthor={CS5486 revision materials}}
\pagestyle{fancy}\fancyhf{}
\fancyhead[L]{\small CS5486｜配套复习题与解析}
\fancyhead[R]{\small 新课件实际页码 1--187}
\fancyfoot[C]{\small\thepage\ / \pageref{LastPage}}
\setlength{\headheight}{14pt}
\setlength{\parskip}{4pt}
\setlist{nosep,leftmargin=1.7em}
\linespread{1.08}
\newcommand{\problem}[3]{\clearpage\section{#1\texorpdfstring{\quad}{ }#2}\label{#1}\noindent\textcolor{compgreen}{\textbf{查表定位：}#3}\par}
\newcommand{\qtext}[1]{\par\textcolor{defblue}{\textbf{复习题题目}}\par #1\par\medskip\hrule\medskip}
\newcommand{\route}[1]{\textbf{圈词与选法：}#1\par}
\newcommand{\given}[1]{\textbf{符号对应：}#1\par}
\newcommand{\solve}{\par\textcolor{calcorange}{\textbf{详细解析}}\par}
\newcommand{\qcheck}[1]{\par\textbf{检查：}#1\par}
\newcommand{\pitfall}[1]{\textcolor{warnred}{\textbf{易错：}}#1\par}
\newcommand{\transfer}[1]{\textbf{换题迁移：}#1\par}

\begin{document}
{\LARGE\bfseries CS5486 期中 Cheatsheet\par 配套复习题与详细解析}\par\medskip
{\large 新课件实际第1--187页｜56道自拟复习题｜中文讲解}\par
\textbf{怎么用这本册子：}先看题目并遮住分隔线后的解析，圈出标签类型、预测输出、损失和约束。到两页cheatsheet所列的“面／栏／编号”找定义或步骤，再自己写出中间结果。做完用“答卷答案、检查、易错、换题迁移”核对；最后完成Q52--Q56混合自测。

\textbf{范围依据：}CS5486a-2026.pdf，共187页。179--182页为SVR，183--185页为LS-SVM回归推导，186页为分类应用，187页为短期预测。开头整学期大纲、人物与应用仅按实际概览深度处理；未展开的后续算法不进入题库。完整逐页覆盖表随项目提供。

\textbf{重点来源与限制：}用户提供的老师消息提到LS-SVM、构建函数与求导、微分线代、方程组、模型共性区别；“最高／最大”所指未明确。本册通过目标极值与最大间隔练习覆盖这些能力，这是复习设计。所有题均为自拟复习题，不是真题，也不表示必考、命题概率、分值或官方评分点。

\textbf{统一约定：}输入按列向量；神经网络目标通常记$z$，输出$a$或$\hat y$；SVM分类标签$y_i\in\{-1,+1\}$；SVR与LS-SVM目标$y_i\in\R$。偏置采用$+b$、阈值$b=-\theta$；$H(0)=1$除非题设另说；BP用$\delta=-\partial E/\partial u$。LS-SVM未知顺序$[\alpha;b]$。每题再列局部具体约定。

\textbf{阅读导航：}
\begin{center}\small\begin{tabular}{p{.44\linewidth}p{.44\linewidth}}
\toprule 正面五栏 & 背面五栏\\\midrule
1：F01/F02概念，M01--M03数学 & 1：S01间隔，S02对偶/KKT\\
2：N01/N02逻辑与状态，D01，P01/P02 & 2：S03软间隔，K01核，D05\\
3：A01 LMS，D02，B01前向 & 3：V01/V02 SVR与偏置\\
4：B02 BP，B03动量/泛化 & 4：L01/L02 LS-SVM建模/系统\\
5：R01/R02，D03/D04比较 & 5：L03例题/条件，D06与关键词索引\\\bottomrule
\end{tabular}\end{center}

\textbf{短答组织：}先完整句写共同任务／结构，再选目标、输出、训练方法等明确维度写差异，最后补适用条件与限制。\textbf{计算组织：}认量→写模型→必要时写目标约束→求导／最优条件→代入消元→回代→解释。

\clearpage\tableofcontents
\problem{Q01}{智能行为、演绎与归纳：怎样写概念短答}{正面第1栏 F01；相关新课件实际第7--12、25、29--35页}
\qtext{以下均为自拟复习题。一个分拣系统有两种工作方式：甲按预设规则“损伤超过阈值则剔除”处理水果；乙根据带标签的水果样本调整模型，再判断新水果。\begin{enumerate}
\item 说明智能系统的含义，并区分甲所用的演绎与乙所用的归纳。
\item 比较 Artificial Intelligence (AI) 与 Computational Intelligence (CI)，以及 hard computing 与 soft computing。评价“AI 不学习，只有 CI 才学习”是否准确。
\item 将 McCulloch--Pitts、Rosenblatt、Widrow、Werbos 与本课模型／方法联系起来；只需人物--贡献，不要求争议年代。
\end{enumerate}}
\route{圈出“预设规则”“从样本”“新水果”“AI/CI”“区别”。查 F01 的定义和比较；本题没有数值目标或更新规则，不应套感知机、BP 或最优化计算式。}
\given{规则与当前事实是甲的已知；样本及其标签是乙的经验；新水果是未见输入。这里比较的是获得／使用知识的方式，不是比较水果特征的数值。}
\solve
\begin{enumerate}
\item 智能系统是表现出某些智能行为的系统，例如推理、从经验学习、适应环境、规划或设计；不要求同时具备全部人类能力。甲从规则和当前事实推出处理结论，是\kw{演绎 deduction}。乙从有限实例形成可用于其他实例的关系，是\kw{归纳 induction}，其推广结论仍可能出错。学习后利用固定参数算新输入，另称预测／推理；此处 inference 不必是形式逻辑证明。
\item AI 是研究和实现智能能力的广泛领域。本课 CI 重点讨论 neural networks、fuzzy logic、evolutionary computation，分别强调学习、对不精确信息的表示及优化等互补能力。课件的传统 AI--CI 对照强调符号／显式规则与适应性取向，但它们并非现代意义下永不交叠的集合；神经网络学习也属于广义 AI。
\item Hard computing 的典型取向是明确模型、精确条件与清晰规则；soft computing 容纳不确定性、不精确性和近似解。“软”不表示没有严格数学，也不保证任何软计算方法都优于精确方法。应依任务、数据和资源选择。
\end{enumerate}
\begin{center}\begin{tabular}{ll}\toprule
人物 & 本课联系\\\midrule
McCulloch--Pitts & 二值神经元与阈值逻辑\\
Frank Rosenblatt & 感知机 Perceptron\\
Bernard Widrow & ADALINE 与误差最小化学习\\
Paul Werbos & 反向传播相关早期工作\\\bottomrule
\end{tabular}\end{center}
\ans{演绎依据给定规则推出结论；归纳从实例建立可推广关系，但需新数据检验。CI 是本课强调神经、模糊、进化方法的家族，能够与其他 AI 方法结合。AI 与学习并不互斥；硬／软计算描述方法取向，不能用来代替具体任务评价。}
\qcheck{逐句检查：是否分别写清“规则\(\to\)结论”和“样本\(\to\)一般关系”；是否提到归纳的不确定性；是否避免把 AI/CI 或 hard/soft 写成绝对互斥。人物表应与模型名称对应。}
\pitfall{用“都是人工智能”代替明确比较；把一个下棋系统的成功写成通用智能证明；把导论中出现的模糊、进化或搜索名称当成本次算法计算范围。}
\transfer{若乙只把原样本存成查找表，要另检验其对未见输入的表现。若系统同时执行规则并训练网络，两种工作方式可以共存，不必二选一。}

\problem{Q02}{ANN 统一框架：结构、五条规则与分类轴}{正面第1栏 F02；正面第2栏 N02；新课件实际第36--50页}
\qtext{某网络有 \(d=2\) 个外部输入、\(N=3\) 个状态节点、\(M=2\) 条可调连接和 \(m=1\) 个外部输出。两个输入节点传值，第三节点聚合加权输入，再经确定性非线性函数；输出选第三节点状态。训练提供目标标签，逐样本调整权重；预测时权重固定，运行连接无环。解释 \(\mathrm{ANN}=(\mathrm{ARCH},\mathrm{RULE})\)，列出有关变量及五条规则，并按课件四个分类轴判断该网络。最后说明训练、推理、状态变化、权重学习的区别。}
\route{圈“变量／规则”“有标签”“确定性”“无环”“固定权重”。查 F02；若问反馈与状态再看 N02。不能由“非线性”推成循环，也不能由“逐样本更新”推成随机激活。}
\given{\(\bm x\in\R^2\) 是外部输入，\(\bm u,\bm v\in\R^3\) 是净输入和激活状态，\(\bm w\in\R^2\) 收集两条连接权重，\(y\in\R\) 是读出。\(z\) 为训练目标，不是预测时必需的输入。}
\solve
ARCH 指部件、连接以及输入输出的组织；RULE 指它们如何运算和学习。只抄五个结构字母仍不足以复现网络，还要说明连接方向、权重及函数。
\begin{center}\begin{tabular}{p{.12\linewidth}p{.77\linewidth}}\toprule
规则 & 本题可对应的实际操作\\\midrule
\(G\) 聚合 & 从输入／旧状态与权重形成净输入，如 \(u_3=w_1x_1+w_2x_2\)。\\
\(F\) 激活 & \(v_3=f(u_3)\)，输入节点直接传值。\\
\(H\) 读出 & \(y=v_3\)，从内部状态取得外部输出。\\
\(E\) 评价 & 例如在给定标签时计算 \(E=\tfrac12(z-y)^2\)。\\
\(L\) 学习 & 根据具体学习规则产生 \(\Delta\bm w\)，令 \(\bm w_{k+1}=\bm w_k+\Delta\bm w_k\)。\\\bottomrule
\end{tabular}\end{center}
实际运行顺序通常是 \(G\to F\to H\)；训练时再评价并更新。\(\mathrm{RULE}=(E,F,G,H,L)\) 是组成清单，非必须按字母顺序执行。

该网络在四轴上是：\kw{确定性}（固定条件给出同一激活）、\kw{前馈}（运行连接无环）、\kw{半线性聚合}（加权和再激活，整体仍可非线性）、\kw{监督学习}（有目标标签）。若聚合含 \(x_1x_2\) 等交互项，才属于本课所说高阶聚合；“隐藏层多”是另一维度。

训练利用数据改变参数；推理以现有参数计算输出。输入变化会使 \(\bm u,\bm v,y\) 变化，并不自动改变 \(\bm w\)。课件的快状态／慢权重对应两种时间尺度，不等于任何前馈网络都有历史记忆。
\ans{ANN 由结构与规则共同定义；本题变量维度分别为 \(2,3,3,2,1\)（按 \(x,u,v,w,y\)）。五规则为聚合、激活、读出、评价和学习。本题是确定性、前馈、加权和型半线性聚合、监督网络；预测中的状态变化与训练中的权重变化是不同动作。}
\qcheck{\(G\) 输出应与 \(u\) 同维，\(F\) 输出与 \(v\) 同维，\(H\) 输出与 \(y\) 同维；权重更新前后维度必须一致。四轴分别检查，不把四个名称当四选一。}
\pitfall{把 \(u\) 与 \(v\) 永远当同一量；以为课件 \(E,L\) 简写没列标签就不需要目标；把实数输入等同连续时间，或把误差用于训练的反馈箭头当作运行连接的循环。}
\transfer{若只去掉标签，监督性必须重判；若只加入状态反馈，前馈／循环必须重判。改变权重初始化方式不必改变前向激活是否确定。连续时间式 \(d\bm w/dt=L\) 给的是变化率，小步更新还需乘时间步长。}

\problem{Q03}{从任务选择输出、监督信息与模型层次}{正面第1栏 F01/F02；正面第2栏 N02；模型细节另见相应栏}
\qtext{对下列任务分别写出需要的输出、是否提供监督信息，以及应先考虑什么模型性质。A：给带“合格／不合格”标签的水果学习分类；B：给实测下一小时负荷学习预测；C：无标签地整理相似观测；D：相同输入脉冲下，输出须取决于先前状态；E：规定输入三个二值开关恰为 \((1,0,1)\) 时报警。评价“只要训练误差最低，所有任务都选同一模型”及“BP 与 MLP 是两种并列网络”两句话。无需替范围外算法设计训练过程。}
\route{圈“标签”“实测数值”“无标签”“先前状态”“规定模式”。先用 F01/F02 区分任务和信息，再查模型。分类／回归由目标输出决定，不能仅凭题中出现“预测”二字判断。}
\given{A 的目标为类别；B 的 \(y\) 是连续实测量；C 只有输入 \(\bm x\)，没有本任务正确类别；D 还需内部状态 \(s_t\)；E 的逻辑要求完全给定。没有给定真实数据分布，不能声称唯一最优模型。}
\solve
\begin{enumerate}
\item A 是监督分类，需要类别或判别分数再决策。若原特征线性可分，可先考虑感知机／线性分类；边界复杂时应考虑非线性表示。需用留出数据评价推广表现。
\item B 是监督回归，输出实数。可比较范围内 RBF、SVR 或 LS-SVM 等回归形式；选择还取决于损失、正则化、数据规模及验证结果。“预测”不自动等于给类别。
\item C 是无监督的结构发现任务，输出可以是组别或简化表示。无监督仍需相似性标准和评价；组号不自动对应“合格类别”。本题只识别任务，不展开后续聚类算法。
\item D 需要历史信息。可用状态递推 \(s_{t+1}=T(s_t,x_t)\) 描述循环结构，必须规定初态和更新方式；普通无状态前馈输入若不含历史，无法区分相同当前输入的不同过去。若显式把历史窗口加入输入，也可用前馈映射处理，需说清输入已改变。
\item E 适合直接构造 M-P／阈值逻辑 decoder。因为目标规则已明确，不必先收集样本让网络学习同一个真值表；N01 给出可核对构造。
\end{enumerate}
模型选择至少要分清三个问题：\kw{表示能力}决定模型能否表达所需关系；\kw{优化}决定训练能否找到较好参数；\kw{泛化}决定未见数据上能否有效。训练误差最低不足以回答后两者之间的权衡，也可能反映过拟合。

MLP 是多层感知机模型／结构；BP 是计算其参数梯度的反向传播方法，随后可由梯度下降等规则更新参数。二者所在层次不同，不能当同类模型二选一。
\ans{A：监督分类；B：监督回归；C：无监督结构发现；D：带状态的映射或显式历史特征；E：已知规则的逻辑构造。选模型应同时检查任务、输入信息、表示、损失与验证表现。BP 是训练中的梯度计算方法，MLP 是模型。}
\qcheck{把每项写成“输入信息\(\to\)输出对象”，检查输出究竟是类别、实数、组别、状态还是逻辑值。对任何推荐加上所需条件，不能把某名称写成无条件最佳。}
\pitfall{“无标签＝没有目标函数”“有反馈＝正在学习”“回归输出再套分类阈值”“MLP与BP并列比较”。范围内的模型选择题可以提到无监督用途，但无需补入本课件未展开的后续算法。}
\transfer{若 B 改成预测“高负荷／低负荷”，应重定义标签与损失；若 A 的输入已含过去十步观测，不能仅凭时间数据就断言模型必须有循环连接。}

\problem{Q04}{向量维度、加权和与仿射映射}{正面第1栏 M01/F02；正面第2栏 N01}
\qtext{给定列向量与参数
\[\bm x=\begin{bmatrix}2\\-1\end{bmatrix},\quad
W=\begin{bmatrix}1&2\\-1&3\end{bmatrix},\quad
\bm b=\begin{bmatrix}1\\-2\end{bmatrix}.
\]
先算 \(\bm u=W\bm x+\bm b\)，再逐分量用 \(H(a)=1\)（\(a\ge0\)）、否则0，得 \(\bm v\)。读出 \(y=[1\ 0]\bm v\)，训练目标 \(z=0\)，评价 \(E=\tfrac12(z-y)^2\)。写出各量维度和完整前向计算。说明 \(W\bm x\)、\(W\bm x+\bm b\)、整个带阈值网络各属于线性、仿射还是非线性。}
\route{圈“列向量”“矩阵乘法”“bias”“逐分量激活”“读出”。查 M01 的维度和仿射定义、F02 的 G/F/H/E；阈值等号查 N01。题目未要求更新权重，不使用学习规则。}
\given{\(\bm x,\bm b,\bm u,\bm v\in\R^2\)，\(W\in\R^{2\times2}\)，读出行向量在 \(\R^{1\times2}\)，\(y,z,E\in\R\)。矩阵的每一行与同一个输入作内积。}
\solve
\begin{enumerate}
\item 聚合：\((2\times2)(2\times1)\) 得 \(2\times1\)，可与同形状 \(\bm b\) 相加。
\[
W\bm x=\begin{bmatrix}1(2)+2(-1)\\-1(2)+3(-1)\end{bmatrix}
=\begin{bmatrix}0\\-5\end{bmatrix},\qquad
\bm u=\begin{bmatrix}0+1\\-5-2\end{bmatrix}=\begin{bmatrix}1\\-7\end{bmatrix}.
\]
\item 激活：\(u_1=1\ge0\)，故 \(v_1=1\)；\(u_2=-7<0\)，故 \(v_2=0\)。于是 \(\bm v=(1,0)\T\)。
\item 读出与评价：\(y=1v_1+0v_2=1\)，\(e=z-y=-1\)，所以 \(E=\tfrac12(-1)^2=\tfrac12\)。误差 \(e\) 与非负代价 \(E\) 是不同量。
\item \(\bm x\mapsto W\bm x\) 为线性变换；加 \(\bm b\ne0\) 后为仿射变换。仿射映射不必过原点；本题在 \(\bm x=\bm0\) 时输出净输入 \(\bm b\ne0\)，已不符合线性变换必须保持零向量的条件。
\item 阶跃激活为非线性，因此整个输入到输出映射一般非线性；不能因聚合是加权和就断言整个网络线性。
\end{enumerate}
\ans{\(W\bm x=(0,-5)\T\)，\(\bm u=(1,-7)\T\)，\(\bm v=(1,0)\T\)，\(y=1,E=1/2\)。前两种映射分别为线性与仿射，加入阈值激活后的整体映射一般非线性。}
\qcheck{用行内积独立重算两个净输入；检查“行数决定输出维数”。所有比较在加完偏置之后执行。若误把 \(W\T\bm x\) 当本题聚合，将得到不同结果，不能只因形状恰好相容就交换。}
\pitfall{把矩阵乘法当逐元素相乘；漏加负偏置；把 \(z-y\) 与其平方误差混同；把非零偏置的映射直接称为线性。这里一般负权按乘法处理，没有 M-P veto。}
\transfer{若题设阈值为 \(\theta_j\)，先比较 \(u_j-\theta_j\)，或把它并入偏置；若题设输出为净输入而没有激活，不应擅自加阶跃。加偏置可写增广输入 \(\tilde{\bm x}=(\bm x\T,1)\T\)。}

\problem{Q05}{链式法则、平方误差梯度与驻点性质}{正面第1栏 M02/M01}
\qtext{设
\[A=\begin{bmatrix}1&2\\0&1\end{bmatrix},\quad
\bm z=\begin{bmatrix}1\\-1\end{bmatrix},\quad
J(\bm w)=\tfrac12\norm{A\bm w-\bm z}^{2},\quad \bm w\in\R^2.
\]
用链式法则求梯度；计算 \(\bm w=(0,0)\T\) 时的梯度；求驻点并判断是最大、最小还是不能判断。另写出 \(\tfrac12\bm w\T Q\bm w\) 的一般梯度，并评价“所有梯度为零的点都是最小值”。}
\route{圈“平方误差”“求导”“驻点”“最大／最小”“二次型”。查 M02，先确定目标方向及变量维度。本题是目标极值判断的复习设计，不推断“最高／最大”具体指什么，也不把所有模型都当最大化。}
\given{令 \(\bm r=A\bm w-\bm z\in\R^2\) 为预测减目标的残差；\(A\in\R^{2\times2}\)。本题 \(J=\tfrac12(r_1^2+r_2^2)\)，前面的 \(1/2\) 会抵消平方求导产生的2。}
\solve
把复合计算拆成 \(\bm w\to\bm r\to J\)。\(\partial J/\partial r_i=r_i\)，而 \(\partial r_i/\partial w_j=A_{ij}\)。对影响同一个 \(w_j\) 的全部残差贡献求和：
\[
\frac{\partial J}{\partial w_j}=\sum_i A_{ij}r_i,
\qquad \nabla J=A\T(A\bm w-\bm z).
\]
明确展开，可独立检查：
\[
J=\tfrac12(w_1+2w_2-1)^2+\tfrac12(w_2+1)^2,
\quad \nabla J=\begin{bmatrix}w_1+2w_2-1\\2w_1+5w_2-1\end{bmatrix}.
\]
在 \(\bm w=\bm0\)，\(\bm r=(-1,1)\T\)，梯度为 \((-1,-1)\T\)。驻点满足
\[
w_1+2w_2=1,\quad 2w_1+5w_2=1.
\]
第二式减第一式的2倍得 \(w_2=-1\)，再回代得 \(w_1=3\)。此时 \(A\bm w=\bm z\)，所以 \(J=0\)。由于平方和非负，零值已经证明它是全局最小值；且 \(A\) 可逆，残差全零解唯一。

二次型需注意矩阵对称性：
\[
\nabla_{\bm w}\!\left(\tfrac12\bm w\T Q\bm w\right)
=\tfrac12(Q+Q\T)\bm w.
\]
仅当 \(Q=Q\T\) 时才能简化为 \(Q\bm w\)。本题 Hessian 为
\(A\T A=\begin{bmatrix}1&2\\2&5\end{bmatrix}\)，对非零 \(\bm h\)，\(\bm h\T A\T A\bm h=\norm{A\bm h}^2>0\)，也验证严格凸性。
\ans{\(\nabla J=A\T(A\bm w-\bm z)\)；原点梯度为 \((-1,-1)\T\)；唯一驻点 \((3,-1)\T\) 是全局最小点，最小值0。驻点本身不自动说明极值性质：\(t^2\) 在0最小，\(-t^2\) 在0最大，\(t^3\) 在0无局部极值。}
\qcheck{将 \((3,-1)\T\) 代回两条梯度式以及 \(A\bm w=\bm z\)，均应成立。维度 \((2\times2)(2\times1)=2\times1\)；若写成 \(A\bm r\)，形状在此例虽相容，链式方向仍错。}
\pitfall{把残差正负与下降方向混淆；漏掉转置；不检查 \(Q\) 对称就写 \(Q\bm w\)；只由 \(\nabla J=0\) 就宣称全局最优。}
\transfer{若损失去掉 \(1/2\)，梯度整体乘2；若加 \(\lambda\norm{\bm w}^2/2\)，梯度加 \(\lambda\bm w\)。若 \(A\) 列不独立，平方和仍为凸目标，但最小点未必唯一。}

\problem{Q06}{增广矩阵、消元与三种解的状态}{正面第1栏 M03/M01}
\qtext{分别把以下系统写成 \(A\bm q=\bm c\)，用增广矩阵行变换求解，并判断唯一解、无解或无穷多解；不使用 \(A^{-1}\bm c\) 作为唯一解法。
\[
\mathrm{I}:\ \begin{cases}2x+y=5,\\x-y=1;\end{cases}\qquad
\mathrm{II}:\ \begin{cases}x+y=2,\\2x+2y=5;\end{cases}\qquad
\mathrm{III}:\ \begin{cases}x+y=2,\\2x+2y=4.\end{cases}
\]}
\route{圈“多元一次”“矩阵”“消元”“判断解”。查 M03：先按同一未知量顺序排列列，再做不改变解集的行操作；看到方程数等于未知数，仍须检查独立性与一致性。}
\given{三个系统均取 \(\bm q=(x,y)\T\in\R^2\)。系数按 \(x,y\) 两列排列，最后一列是右端。\(R_i\) 表示第 \(i\) 行，增广分隔线不代表额外未知量。}
\solve
\kw{I：}\(A=\begin{bmatrix}2&1\\1&-1\end{bmatrix}\)，\(\bm c=(5,1)\T\)。
\[
\left[\begin{array}{cc|c}2&1&5\\1&-1&1\end{array}\right]
\xrightarrow{R_1\leftrightarrow R_2}
\left[\begin{array}{cc|c}1&-1&1\\2&1&5\end{array}\right]
\xrightarrow{R_2\leftarrow R_2-2R_1}
\left[\begin{array}{cc|c}1&-1&1\\0&3&3\end{array}\right].
\]
第二行给 \(3y=3\)，所以 \(y=1\)；回代 \(x-y=1\) 得 \(x=2\)。两列均有主元，\(\operatorname{rank}A=\operatorname{rank}[A|\bm c]=2\)，故唯一解。

\kw{II：}\(A=\begin{bmatrix}1&1\\2&2\end{bmatrix}\)，\(\bm c=(2,5)\T\)。
\[
\left[\begin{array}{cc|c}1&1&2\\2&2&5\end{array}\right]
\xrightarrow{R_2\leftarrow R_2-2R_1}
\left[\begin{array}{cc|c}1&1&2\\0&0&1\end{array}\right].
\]
第二行要求 \(0x+0y=1\)，矛盾。\(\operatorname{rank}A=1\)，增广矩阵秩2，故无解；不应继续把0当除数。

\kw{III：}\(A\) 与II相同，\(\bm c=(2,4)\T\)。
\[
\left[\begin{array}{cc|c}1&1&2\\2&2&4\end{array}\right]
\xrightarrow{R_2\leftarrow R_2-2R_1}
\left[\begin{array}{cc|c}1&1&2\\0&0&0\end{array}\right].
\]
第二行是恒等式，只有一个独立约束。令自由变量 \(y=t\)，则 \(x=2-t\)，\(t\in\R\)。两矩阵秩同为1，小于未知数2，故无穷多解。
\ans{I：\((x,y)=(2,1)\)，唯一解。II：出现矛盾行 \([0\ 0|1]\)，无解。III：\((x,y)=(2-t,t)\)，\(t\in\R\)，无穷多解。一般先查一致性；一致且主元数等于未知数则唯一，一致但有自由变量则无穷多。}
\qcheck{I 回代得 \(2(2)+1=5\)、\(2-1=1\)；III 回代得 \((2-t)+t=2\) 及其2倍4。II 原式第二左端必须是第一左端2倍，右端却不是，直接暴露矛盾。}
\pitfall{只改系数不改右端；把 \([0\ 0|0]\) 与 \([0\ 0|1]\) 混同；认为方程数等于未知数就一定唯一；列交换后忘记同步交换未知量解释。}
\transfer{把II右端5改成4，矛盾消失但没有增加独立约束，因此变成无穷多解。后面 LS-SVM 的分块系统同样按已声明未知量列顺序消元和回代。}

\problem{Q07}{M-P 的 veto、逻辑门与阈值边界}{正面第2栏 N01；新课件实际第51--56页}
\qtext{所有输入均为 \(\{0,1\}\)，达到阈值输出1。\begin{enumerate}
\item M-P 单元的前两个输入兴奋、第三个输入抑制，阈值 \(\theta=1\)。求 \(\bm x=(1,1,1)\) 的输出；再与普通加权 TLU 的 \(\bm w=(1,1,-1)\)、\(\theta=1\) 比较。
\item 用 M-P 单元设计两输入 AND、OR 和一输入 NOT，给出连接性质与阈值，并核对真值表。
\end{enumerate}}
\route{圈“M-P”“抑制”“veto”“普通加权”“达到阈值”。查 N01，先判题目是哪种负连接语义，再算；不能将二者混成一种规则。}
\given{M-P 中 active inhibitory input 强制输出0；仅在无抑制输入激活时才比较兴奋输入个数与 \(\theta\)。普通 TLU 使用 \(s=\bm w\T\bm x\)，输出 \(H(s-\theta)\)，没有额外否决。}
\solve
\begin{enumerate}
\item M-P：\(x_3=1\) 是激活的抑制输入，故立即 \(y=0\)，不管另两个兴奋输入有多强。普通 TLU：\(s=1+1-1=1\ge\theta\)，故 \(y=1\)。负权贡献能被正贡献抵消，veto 不能。
\item AND：两输入均兴奋、\(\theta=2\)；OR：两输入均兴奋、\(\theta=1\)。它们的兴奋和均为 \(x_1+x_2\)，只因阈值不同产生不同功能。
\end{enumerate}
\begin{center}\begin{tabular}{ccccc}\toprule
\(x_1\)&\(x_2\)&兴奋和&AND \(\theta=2\)&OR \(\theta=1\)\\\midrule
0&0&0&0&0\\0&1&1&0&1\\1&0&1&0&1\\1&1&2&1&1\\\bottomrule
\end{tabular}\end{center}
NOT 取一个抑制输入、\(\theta=0\)：\(x=1\) 时 veto 输出0；\(x=0\) 时没有抑制，兴奋和为0，且 \(0\ge0\)，所以输出1。普通 TLU 也可用 \(w=-1,\theta=0\) 实现这个特定 NOT，但这不使所有普通负权变成 veto。
\begin{center}\begin{tabular}{ccc}\toprule
\(x\)&是否 veto&NOT输出\\\midrule 0&否，\(0\ge0\)&1\\1&是&0\\\bottomrule
\end{tabular}\end{center}
\ans{给定输入下 M-P 输出0，普通 TLU 输出1。AND：两兴奋输入、阈值2；OR：两兴奋输入、阈值1；NOT：一个抑制输入、阈值0。所有边界按“\(\ge\)”触发，且 veto 优先。}
\qcheck{对每个门核对全部输入而非只试一行。特别检查 OR 的和1、AND 的和2和 NOT 的无输入兴奋和0，均恰在边界，应按约定输出1。}
\pitfall{先求带 \(-1\) 的总和再处理 M-P，导致 veto 被抵消；把“严格大于”误当课件“不小于”；觉得无兴奋输入就绝不能输出1而算错 NOT。}
\transfer{\(n\) 输入 AND 取 \(n\) 个兴奋输入与 \(\theta=n\)，OR 仍为 \(\theta=1\)。若普通 TLU 的权重与阈值同乘正数，预测不变；负数缩放会改变不等号方向，须重新判断。}

\problem{Q08}{Decoder：从目标模式设计并核对全部输入}{正面第2栏 N01；新课件实际第57--61页}
\qtext{构造一个三输入 decoder，仅在 \(\bm r=(1,0,1)\) 时输出1。分别给出 M-P veto 实现和普通加权 TLU 实现；列出完整八行真值表。说明普通 TLU 对一般目标 \(\bm r\in\{0,1\}^n\) 的构造，以及为什么只识别一个模式。}
\route{圈“仅在”“完整模式”“decoder”。查 N01；本题规则已知，做参数构造与真值表验证，不需要感知机训练。先区分 veto 与可相抵的普通负权。}
\given{输入均0/1，输出规则达到阈值触发。目标1的位为第1、3位，目标0的位为第2位；普通 TLU 的分数记 \(s\)，避免与激活后的0/1输出混淆。}
\solve
M-P：第1、3位兴奋，第2位抑制，\(\theta=2\)。若 \(x_2=1\) 立即否决；若 \(x_2=0\)，仍须 \(x_1+x_3\ge2\)，所以必须 \(x_1=x_3=1\)。这正好是指定模式。

普通 TLU：取 \(\bm w=(1,-1,1)\T\)、\(\theta=2\)，即 \(y=H(x_1-x_2+x_3-2)\)。完整核对如下；两种构造的最终输出相同。
\begin{center}\begin{tabular}{ccccc}\toprule
\(x_1\)&\(x_2\)&\(x_3\)&普通分数 \(s=x_1-x_2+x_3\)&两实现输出\\\midrule
0&0&0&0&0\\0&0&1&1&0\\0&1&0&$-1$&0\\0&1&1&0&0\\
1&0&0&1&0\\1&0&1&2&1\\1&1&0&0&0\\1&1&1&1&0\\\bottomrule
\end{tabular}\end{center}
一般构造令 \(k=\sum_i r_i\) 为目标中1的个数，取
\[
w_i=2r_i-1,\qquad \theta=k,\qquad y=H(\bm w\T\bm x-k).
\]
完全匹配时，目标1的位各贡献1，目标0的位贡献0，故总分 \(k\)。目标1误成0，贡献少1；目标0误成1，贡献变成 \(-1\)，也少1。因此每错一位总分就从 \(k\) 减1，只要有错位便低于阈值。这是该简单构造的教学推导。
\ans{M-P：\(x_1,x_3\) 兴奋、\(x_2\) 抑制、\(\theta=2\)。普通 TLU：\(\bm w=(1,-1,1)\T,\theta=2\)。两者均只在101输出1。一般目标可取 \(w_i=2r_i-1,\theta=\sum_i r_i\)，任一错位使分数降低。}
\qcheck{至少检查目标101、含激活抑制的111、少一个所需兴奋的001；完整表确认没有遗漏其他触发输入。参数不必唯一，真正验收条件是全部行正确。}
\pitfall{目标为0的位直接不连接，会把“该位必须0”误成“不关心该位”；把单次匹配识别理解成某种学习收敛；沿用旧课件页码找错 decoder。}
\transfer{若目标改成001，普通参数为 \((-1,-1,1)\T,\theta=1\)；若全零目标，则 \(k=0\)，全负权、阈值0恰好只接收全零。多个正模式可各设 decoder，再用 OR 合并。}

\problem{Q09}{多层 XOR：表示能力与单元限制}{正面第2栏 N01/P02；正面第3栏 B01}
\qtext{输入 \((x_1,x_2)\in\{0,1\}^2\)，XOR 在01、10输出1，在00、11输出0。用两层 M-P 计算单元构造 XOR，写出隐藏量及真值表。再证明任意普通单个加权阈值单元也不能实现 XOR。最后说明“可以构造”是否意味着“某训练算法必能学到且泛化良好”。}
\route{圈“XOR”“两层”“单个阈值”“不能”。查 N01 的正行 decoder→OR，P02 的线性可分限制及 B01 的隐藏层表示；不要反复调单个单元的学习率来寻找本不存在的解。}
\given{所有激活取达到阈值输出1，输入层只传值，不算计算层。第一层设两个模式 detector \(h_1,h_2\)，第二层输出 \(y\)。单个 TLU 使用任意实数 \(w_1,w_2,\theta\)，并非只允许正权。}
\solve
取
\[
h_1=x_1\land\neg x_2,\qquad
h_2=\neg x_1\land x_2,\qquad y=h_1\lor h_2.
\]
\(h_1\)：\(x_1\) 兴奋、\(x_2\) 抑制，\(\theta=1\)；\(h_2\) 对调兴奋与抑制，\(\theta=1\)；输出单元两个兴奋输入、\(\theta=1\)。因此总共三个计算单元、两层计算，隐藏层把两个不同模式分别识别。
\begin{center}\begin{tabular}{ccccc}\toprule
\(x_1\)&\(x_2\)&\(h_1\)&\(h_2\)&\(y=h_1\lor h_2\)\\\midrule
0&0&0&0&0\\0&1&0&1&1\\1&0&1&0&1\\1&1&0&0&0\\\bottomrule
\end{tabular}\end{center}
若单个普通 TLU 可以实现它，四行要求
\[
0<\theta,\quad w_2\ge\theta,\quad w_1\ge\theta,
\quad w_1+w_2<\theta.
\]
中间两式相加给 \(w_1+w_2\ge2\theta>\theta\)，与最后一式矛盾。因此偏置 \(b=-\theta\) 或允许负权也不能解除单个线性阈值单元的限制。多层构造增加的是表示能力。

此构造给出一套能实现真值表的参数，尚未证明某种优化算法能自动找到它；有限四行表示正确也未提供任意实际数据的泛化保证。表示、优化、泛化三个问题须分别论证。
\ans{用两个带否定的 AND detector 加一个 OR 即实现 XOR；单个 TLU 的四个阈值约束彼此矛盾，所以无法实现。两层网络存在可用参数，不等于任意初始化、训练法或新数据都能成功。}
\qcheck{逐行检查隐藏输出，尤其11时两隐藏单元都因对方抑制而输出0。检查“不可能证明”没有偷偷限制权重非负，且保留负类严格小于与正类不小于的边界区别。}
\pitfall{将两个原输入直接 OR；把两层计算误写成只有两个计算单元；把 BP 当成与 MLP 并列的模型；将单层 XOR 训练不终止归咎于训练轮数不足。}
\transfer{任意 Boolean 函数都可为每个正行设 decoder 后 OR；这种直接构造可能很大，并不保证最少单元。无抑制兴奋网络只产生单调输出，不能构造 NOT 这样的下降关系。}

\problem{Q10}{循环网络：用旧状态作同步更新}{正面第2栏 N02/N01；新课件实际第62--64页}
\qtext{本题规定同步离散更新：每个单元延迟一步，所有新值均用旧状态。初态 \((A_0,B_0)=(0,0)\)。输入序列为 \(u_0,u_1,u_2,u_3=1,0,1,0\)，状态规则给定为
\[
A_{t+1}=(u_t\lor A_t)\land\neg B_t,
\qquad B_{t+1}=(u_t\land A_t)\land\neg B_t.
\]
每次更新后以 \(B_{t+1}\) 为本次输出。列出状态和输出；说明为什么它有历史依赖，以及有限自动机与循环 M-P 网络的关系。}
\route{圈“初态”“同步”“旧状态”“更新后输出”。查 N02 的状态递推步骤，再查 N01 的逻辑门。原课件图没有完整给出时序，本题已补明，不能改成逐个立即使用新状态。}
\given{\(u_t\in\{0,1\}\) 是外部输入；\((A_t,B_t)\) 是内部状态，均二值。每步先冻结 \((A_t,B_t)\)，分别算两个新值，最后同时写回。这里改变的是状态，权重未学习。}
\solve
\begin{center}\begin{tabular}{ccccc}\toprule
\(t\)&\(u_t\)&旧状态 \((A_t,B_t)\)&新状态 \((A_{t+1},B_{t+1})\)&输出\\\midrule
0&1&(0,0)&(1,0)&0\\
1&0&(1,0)&(1,0)&0\\
2&1&(1,0)&(1,1)&1\\
3&0&(1,1)&(0,0)&0\\\bottomrule
\end{tabular}\end{center}
首步 \(A_1=(1\lor0)\land1=1\)，\(B_1=(1\land0)\land1=0\)。第二步外部输入为0，但 \(A\) 的自反馈让它保持1。第三步 \(u_2=A_2=1\) 且 \(B_2=0\)，所以 \(B_3=1\)。第四步旧 \(B_3=1\) 同时否决两个新状态，清零。

相同当前输入 \(u=1\)，旧状态00产生新状态10、输出0；旧状态10则产生11、输出1。输出的差别由历史保存的状态引起，而非当前输入不同。只有输出位 \(B\) 也不足以代表完整状态，因为00与10的 \(B\) 同为0，但对下一输入反应不同。

有限自动机只有有限状态和有限种输入，下一状态由旧状态和当前输入决定。用有限个 Boolean 位编码状态，再以 M-P 网络实现每个下一状态位的 Boolean 逻辑，配合反馈与适当延迟，即能模拟有限自动机。课件的命题是表示／模拟能力，不是权重自动学习的保证。本例理论有4个二位状态，从00按给定规则可达00、10、11。
\ans{新状态依次10、10、11、00，更新后输出为0、0、1、0。反馈保存历史，当前输入与旧状态共同决定下一步；有限自动机可由循环 M-P 网络模拟，但须规定状态编码与时序。}
\qcheck{每一行的旧状态都等于上一行的新状态；计算 \(B_1\) 时使用 \(A_0=0\)，不能偷用 \(A_1=1\)。最后旧 \(B=1\) 的 veto 应使两新状态均0。}
\pitfall{把状态更新当权重训练；忽略初态；只给输入不规定同步/异步；把有限状态误说成可记住任意长历史的全部细节。}
\transfer{若初态改成10，第一次 \(u=1\) 就会输出1，必须重算整条状态链。若连续输入全1，本规则从00按00→10→11→00循环，不能无条件称为对任意脉冲流精确“除以二”。}

\problem{Q11}{M-P V.S. Perceptron：共同点、差异与限制}{正面第2栏 D01/N01/P01/P02}
\qtext{用完整句子比较 M-P 神经元与简单感知机。答案须包含共同点、连接／激活与学习方面的差异、适用任务和限制。评价三句话：(1)“它们都有阈值，所以负权一定有 veto”；(2)“感知机只要训练足够久，就能在原二维输入上学会 XOR”；(3)“M-P 固定参数不学习，因此不能被组织成循环网络”。}
\route{圈“compare”“共同点”“区别”“限制”。查 D01，分别按模型定义、参数来源、任务、能力比较；不能仅写“都是神经网络”，也不能把一个单元的限制直接扩大为所有多层网络。}
\given{本课 M-P 指0/1、兴奋／绝对抑制与固定阈值逻辑的定义；感知机采用实数 \(\bm w,b\)，\(y=H(\bm w\T\bm x+b)\)。若写学习式，先声明标签为0/1、\(H(0)=1\)。}
\solve
两者都把多项输入连接到计算单元，经阈值产生二值输出；都可放入更大的网络，并以逻辑／几何方式分析表示能力。这是比“都是机器学习”更具体的共性；M-P 本身未规定从样本学习。
\begin{center}\begin{tabular}{>{\raggedright\arraybackslash}p{.16\linewidth}>{\raggedright\arraybackslash}p{.34\linewidth}>{\raggedright\arraybackslash}p{.37\linewidth}}\toprule
比较维度&M-P&简单 Perceptron\\\midrule
参数与连接&按逻辑设计固定兴奋、抑制连接及阈值&实数权重和偏置，可根据带标签样本调整\\
负作用&激活的抑制输入直接 veto&一般负权降低净输入，可被其他正贡献抵消\\
学习规则&该单元定义没有给训练法&误差纠正，例如0/1约定下 \(\Delta\bm w=\eta(z-y)\bm x\)，\(\Delta b=\eta(z-y)\)\\
典型任务&已知逻辑门、decoder、有限状态逻辑构造&监督二分类，学习一个线性阈值边界\\
能力限制&单个单元受其连接和阈值形式限制；可多层组合&单个原空间线性阈值边界，XOR不可分\\\bottomrule
\end{tabular}\end{center}
感知机收敛保证需要有限训练集线性可分及相应更新条件；不保证任意数据都收敛、最大间隔或未见样本全正确。M-P 的多层构造能扩展 Boolean 表示能力，但不等于自带训练算法。

(1)错：阈值激活与 veto 是两个不同规则，普通负权没有一票否决语义。(2)错：XOR 对单个线性阈值单元不可表示，增加轮数不产生可行参数。(3)错：连接结构与是否训练是不同维度；固定权重的循环网络仍能更新内部状态。
\ans{M-P与感知机共同使用连接、聚合和二值阈值输出；前者强调固定逻辑构造和绝对抑制，后者以可学习实数权重／偏置完成监督线性分类。单个感知机不能表示XOR；多层表示、参数学习与循环状态是不同设计维度。三句话均不正确。}
\qcheck{答案是否同时覆盖共性、明确维度的差异、任务与条件；是否将“不能学会XOR”限定为原输入上的单个线性阈值单元；是否把收敛与泛化分开。}
\pitfall{称 M-P 与感知机都必须从数据训练；把普通负权解释成 veto；称感知机训练求唯一解或最大间隔；以为“循环”必然意味着“在线学习”。}
\transfer{若题给普通加权 TLU 也称“M-P”，先按题目明确公式作答，并指出与课件 veto 定义的不同。改成多层网络后，必须重新讨论表示能力，不能沿用单单元XOR结论。}

\problem{Q12}{混合纠错：由文字建立回归式、梯度与解集}{正面第1栏 F02/M01/M02/M03；梯度学习另见正面第3栏 A01}
\qtext{欲用两个实数特征预测一个连续数值，模型含偏置，按全部样本半平方误差之和评价。给两样本 \(\bm x_1=(1,0)\T,z_1=1\)，\(\bm x_2=(0,1)\T,z_2=2\)。一份解答声称：“\(f(\bm x)=\bm w\T\bm x+b\) 即使 \(b\ne0\) 仍是线性变换；损失梯度为 \(\sum_i(z_i-f_i)\bm x_i\)，按负梯度更新；梯度零表示最大值；两条方程能唯一确定 \(w_1,w_2,b\)。”请逐项纠正，写出模型与目标，从零参数、\(\eta=0.1\) 作一次批量梯度更新，并求全部零误差解。}
\route{圈“连续数值”“含偏置”“全部样本半平方误差”“批量”“唯一”。查 M01/M02/M03 与 F02；本题是可微实数回归，不能把最终输出换成感知机类别后再算损失。}
\given{未知为 \(\bm w=(w_1,w_2)\T\)、\(b\)。预测 \(f_i=\bm w\T\bm x_i+b\)，目标为 \(z_i\)。本题用损失\emph{和}，没有再除样本数；批量梯度中的所有残差均由同一套旧参数产生。}
\solve
由文字建立
\[
f(\bm x)=\bm w\T\bm x+b,\qquad
\min_{\bm w,b}J=\tfrac12\sum_{i=1}^2(f_i-z_i)^2
=\tfrac12(w_1+b-1)^2+\tfrac12(w_2+b-2)^2.
\]
有非零偏置时为仿射映射。求导得到
\[
\nabla_{\bm w}J=\sum_i(f_i-z_i)\bm x_i,
\qquad \frac{\partial J}{\partial b}=\sum_i(f_i-z_i).
\]
原解答把梯度与负梯度符号混了；若用 \(z_i-f_i\)，应将其作为下降\emph{更新方向}相加，而不能再取一次负号。

旧参数全零时 \(f_1=f_2=0\)，所以
\[
\nabla_{\bm w}J=(-1,-2)\T,\quad J_b=-3,
\quad \bm w^+=(0.1,0.2)\T,\quad b^+=0.3.
\]
回算新预测 \(f_1^+=0.4,f_2^+=0.5\)，新损失
\[
J^+=\tfrac12[(0.4-1)^2+(0.5-2)^2]
=\tfrac12(0.36+2.25)=1.305,
\]
低于旧损失 \(J=\tfrac12(1+4)=2.5\)。这是给定步长的一次下降检查，不是对任意步长的保证。

全部零误差解满足
\[
\left[\begin{array}{ccc|c}1&0&1&1\\0&1&1&2\end{array}\right],
\qquad w_1+b=1,\quad w_2+b=2.
\]
系数秩2、未知数3，无矛盾但有自由变量。令 \(b=t\)，得 \(w_1=1-t,w_2=2-t\)，\(t\in\R\)。这些解均令非负平方和为0，故均为全局最小点，说明本题最小点并不唯一；“驻点就是最大”也错误。
\ans{正确梯度是预测减目标。一次更新为 \((w_1,w_2,b)=(0.1,0.2,0.3)\)，新损失1.305。全部零误差最小解为 \((1-t,2-t,t)\)，\(t\in\R\)。非零偏置使映射为仿射；驻点、极值性质与唯一性须分别判断。}
\qcheck{重算两新预测和损失；把任意 \(t\) 的解代回两样本，均得1、2。梯度的权重部分是二维，偏置部分是标量；每个样本同时影响偏置，因此偏置梯度是两个残差之和。}
\pitfall{用新 \(w_1\) 计算第二个样本梯度，意外改成在线更新；将损失和换成平均却沿用同一数值梯度；因有零误差解就推断未来样本预测唯一。}
\transfer{如加正则化，目标和驻点都须重写，不能直接沿用该自由解族；如补第三条独立样本约束，可能得到唯一解，也可能冲突，仍须消元检查。}

\problem{Q13}{感知机：按样本顺序学习 OR}{正面第2栏／P01；可分性见 P02}
\qtext{输入按 $(0,0),(0,1),(1,0),(1,1)$ 顺序反复出现，目标依次为 $z=0,1,1,1$。规定 $u=w_1x_1+w_2x_2+b$，$\hat z=H(u)$，其中 $H(u)=1$ 当 $u\ge0$，否则为0。初值 $w_1=w_2=b=0$，学习率 $\eta=1$。用 $\tilde w\leftarrow\tilde w+\eta(z-\hat z)\tilde x$ 更新，$\tilde x=(x_1,x_2,1)\T$。完成第一轮，判断是否收敛，再继续到一整轮没有误分类。}
\route{圈出“\kw{阈值输出、0/1目标、逐样本更新、顺序}”，查 P01。这里用类别差 $z-\hat z$；不能用 ADALINE 的连续残差 $z-u$。P02 提醒：OR 可分，标准感知机训练有有限步收敛保证。}
\given{$x\in\R^2$，$w\in\R^2$，$b\in\R$；增广后 $\tilde x,\tilde w\in\R^3$。阈值形式为 $w\T x-\theta$ 时，$b=-\theta$。每个样本先用当前全部旧参数预测，再更新全部分量。}
\solve
第一轮的每一行都从上一行更新后的参数开始：
\[
\begin{array}{c|r|c|r|c}
 x&u&\hat z&z-\hat z&\text{更新后 }(w_1,w_2,b)\\\hline
 (0,0)&0&1&-1&(0,0,-1)\\
 (0,1)&-1&0&1&(0,1,0)\\
 (1,0)&0&1&0&(0,1,0)\\
 (1,1)&1&1&0&(0,1,0)
\end{array}
\]
例如第二行的代入是
\[
 (0,0,-1)+(1-0)(0,1,1)=(0,1,0).
\]
第一轮结束后，$(0,0)$ 仍给 $u=0$、预测1，因此尚未收敛。继续按同一顺序扫描：
\begin{itemize}
\item 第二轮遇 $(0,0)$ 时更新为 $(0,1,-1)$；遇 $(0,1)$ 无需更新；遇 $(1,0)$ 时 $u=-1$、误分类，更新为 $(1,1,0)$；最后一点正确。
\item 第三轮遇 $(0,0)$ 时更新为 $(1,1,-1)$；其余三点都正确。
\item 第四轮四点全部正确，没有任何更新，才满足停止检查。
\end{itemize}
\ans{第一轮参数为 $(0,1,0)$，尚未收敛；最终可得 $(1,1,-1)$，判别式 $\hat z=H(x_1+x_2-1)$，共发生5次非零更新。}
\qcheck{最终四点分数为 $-1,0,0,1$，按 $H(0)=1$ 得 $0,1,1,1$。更新次数不等于样本访问次数；最后还需扫描一轮确认。}
\pitfall{把 $H(0)$ 当0会改变整条轨迹。漏掉增广输入的末项1，就无法更新 bias；若直接更新阈值，应为 $\theta\leftarrow\theta-\eta(z-\hat z)$。}
\transfer{换样本顺序或初值，路径和最终分隔线可能改变；仍先预测再纠错并检查整轮。若改成 XOR，先查 P02 判断结构是否有解，增加轮次不改变单阈值的表示限制。}

\problem{Q14}{为什么 XOR 不能由一个感知机实现？}{正面第2栏／P02；多层构造见正面第3栏／B01}
\qtext{令 XOR 目标为 $z(0,0)=z(1,1)=0$，$z(0,1)=z(1,0)=1$。单个感知机使用 $\hat z=H(w_1x_1+w_2x_2+b)$，$H(0)=1$。用不等式证明不存在正确分类四点的参数。再解释：“感知机不收敛只是学习率太小”“有限数据可分就能严格分开”“收敛后一定得到最大间隔”三句话。}
\route{圈出“\kw{XOR、single perceptron、linear separability、convergence}”，查 P02。题目先问模型能否表示；不应先套 LMS 或 BP 调学习率。多层网络的非线性中间表示才可能改变可分性。}
\given{$u=w\T x+b$ 是仿射分数；$z=1$ 要求 $u\ge0$，$z=0$ 要求 $u<0$。参数 $w_1,w_2,b$ 是未知量，四个样本给四条约束。}
\solve
逐点写出条件：
\[
\begin{array}{c|c}
 (0,0):z=0&b<0\\
 (0,1):z=1&w_2+b\ge0\\
 (1,0):z=1&w_1+b\ge0\\
 (1,1):z=0&w_1+w_2+b<0
\end{array}
\]
两个正类条件相加，得 $w_1+w_2+2b\ge0$；两个负类条件相加，却得 $w_1+w_2+2b<0$，矛盾。因此参数不存在。

几何上，两个正类顶点连成的线段和两个负类顶点连成的线段共用中点 $(1/2,1/2)$。仿射分数在一条线段中点等于两端分数的平均；同一中点不可能同时属于正、负两侧。

三句判断分别为：
\begin{enumerate}
\item 第一项错误。对单阈值 XOR，失败来自表示能力；调整步长、训练更久不能产生不存在的解。一般数据不收敛也可能有实现或优化原因，不能只凭现象断言是哪一种。
\item 第二项成立，但要保留\kw{有限数据集}条件。给定一侧 $u\ge0$、另一侧 $u<0$，有限个负类分数的最大值 $a$ 仍小于0；把分界水平从0移到 $a/2$，即可使两侧都严格分开。
\item 第三项错误。感知机收敛定理保证在有限、线性可分训练集上找到分隔参数；它没有以间隔为优化目标。最大间隔要查背面 SVM。
\end{enumerate}
\ans{XOR 的四条线性判别约束相互矛盾，因此一个感知机无法实现。感知机收敛须有可分解；该保证既不是最大间隔保证，也不是测试集泛化保证。}
\qcheck{把候选参数逐点代入最可靠；只画一条看似正确的线不算证明。OR 可由 $w=(1,1)\T,b=-1/2$ 严格分开，可用来对照 XOR。}
\pitfall{“线性分类器”允许 bias，因此边界一般不经过原点。证明时必须照标签写不等号，不能把目标0误写成“分数等于0”。}
\transfer{如果换成更高维特征或含非线性隐藏层，先在新表示空间重新写约束；原始输入空间不可分，并不推出任何模型都不可分。纯线性或仿射层叠加仍只有仿射分数。}

\problem{Q15}{输入空间、权空间与容量计数}{正面第2栏／P02}
\qtext{对带 bias 的二维感知机，说明固定权重与固定样本时分别在什么空间得到超平面。将 OR 四点的分类要求写成权空间约束。设 $m$ 个过原点的超平面位于 $D$ 维权空间且处于一般位置，其最大开区域数为 $R(m,D)=2\sum_{k=0}^{D-1}\binom{m-1}{k}$。计算 $R(4,3)$；解释为什么两个二值输入的16个布尔函数中只有14个为阈值函数，并比较三输入的256个布尔函数、104个阈值函数和128区上界。}
\route{圈出“\kw{input/weight space、duality、regions、Boolean/threshold functions}”，查 P02。这里的对偶是交换“固定样本/固定参数”的几何观点；它与 SVM 的 Lagrangian 对偶不是同一推导。}
\given{输入数 $d=2$，$\tilde x=(x_1,x_2,1)\T$，$\tilde w=(w_1,w_2,b)\T$，参数空间维数 $D=d+1=3$。每个固定样本给一个过权空间原点的平面 $\tilde w\T\tilde x=0$。组合数在 $k>m-1$ 时取0。}
\solve
固定 $\tilde w$ 而改变 $x$，得到输入空间中的边界 $w_1x_1+w_2x_2+b=0$。固定 $x$ 而改变参数，得到权空间中的平面；目标标签选择其哪一侧。OR 的可行参数交集为
\[
 b<0,\quad w_2+b\ge0,\quad w_1+b\ge0,\quad w_1+w_2+b\ge0.
\]
例如 $(1,1,-1/2)$ 位于这个交集中。感知机训练是在权空间寻找一个可行点；输入点本身没有被更新。

对三维中心排列，
\[
 R(4,3)=2\left[\binom30+\binom31+\binom32\right]
 =2(1+3+3)=14.
\]
这里的等号要求一般位置：法向量没有额外退化，使每次切割达到最大数；否则实际区域数至多为该数。四平面都过原点的条件不能删掉，不能把该式当任意位置平面的分区公式。

两个二值输入有 $2^2=4$ 行真值表，每行独立选0或1，所以共有 $2^4=16$ 个函数。权空间每个开区域给一种严格符号组合；这四点的增广平面达到14区，缺少的两种就是 XOR 和 XNOR。

对 $d=3$，输入组合数 $m=8$，含偏置参数维数 $D=4$，故
\[
 R(8,4)=2(1+7+21+35)=128<2^8=256.
\]
课件给实际阈值函数数104。立方体顶点的增广法向量不满足所需的一般位置条件，所以上界128不必达到；正确关系为 $104\le128<256$。
\ans{输入空间边界固定参数、权空间边界固定样本。$R(4,3)=14$；布尔函数总数 $2^{2^d}$，阈值函数只占子集，且 $T_d\le R(2^d,d+1)$。}
\qcheck{维数使用 $D=d+1$，故二维输入不能误代 $R(4,2)=8$。三输入的8是样本组合数，256才是布尔函数数，104是课件列出的实际阈值函数数。}
\pitfall{最大区域数、某组样本的实际区域数和全部布尔函数数是三种量。“更多训练轮次”只能寻找已有可行区域，不能增加该结构能表达的标签模式。}
\transfer{换 $m,D$ 后，组合求和步骤不变；先重新判断是否为中心排列、是否一般位置。输入数增加时参数维数也变，不能把“固定 $D$ 时随 $m$ 为 $D-1$ 次多项式”无条件跨维度套用。}

\problem{Q16}{阈值、bias 与单极／双极编码}{正面第2栏／P01、P02}
\qtext{OR 用 $x_i\in\{0,1\}$ 表示，分数为 $u=x_1+x_2-1$，输出 $z=H(u)$，$H(0)=1$。将输入改为 $s_i=2x_i-1\in\{-1,1\}$，目标改为 $t=2z-1$，写出等价分数与输出。再用双极感知机更新 $\tilde v\leftarrow\tilde v+\eta(t-\hat t)\tilde s$：初参数全0，$\eta=1/2$，先看原始输入 $(0,0)$，完成一次更新。最后判断“双极编码一定有更高准确率”。}
\route{圈出“\kw{unipolar/bipolar、threshold、bias、encoding}”，查 P01/P02。先转换输入坐标和目标编码，再套相应误差；不能只把0替换为−1而保留全部旧参数。}
\given{$x_i=(s_i+1)/2$；$\tilde s=(s_1,s_2,1)\T$，$\tilde v=(v_1,v_2,b_s)\T$；双极输出约定为 $\hat t=2H(u)-1$，故 $u=0$ 时输出 $+1$。}
\solve
代入坐标变换：
\[
 u=\frac{s_1+1}{2}+\frac{s_2+1}{2}-1
 =\frac12s_1+\frac12s_2.
\]
因此等价参数是 $(v_1,v_2,b_s)=(1/2,1/2,0)$，输出为 $\hat t=2H((s_1+s_2)/2)-1$。四个原输入转换成 $(-1,-1),(-1,1),(1,-1),(1,1)$，分数依次为 $-1,0,0,1$，目标为 $-1,+1,+1,+1$，保持相同分类。

训练小问另从\kw{全零参数}出发，并非从上述等价参数出发。第一样本为 $\tilde s=(-1,-1,1)\T$、$t=-1$。旧分数0，预测 $\hat t=+1$，所以
\[
 t-\hat t=-2,\qquad
 \tilde v'=(0,0,0)\T+\frac12(-2)(-1,-1,1)\T=(1,1,-1)\T.
\]
新分数为 $1(-1)+1(-1)-1=-3$，第一样本已正确。

“双极编码一定更准”不成立。编码会改变几何参数化、误差步幅和不同分类在权空间中的比例，可能有代数便利；它不会让 XOR 从不可分变可分，也不能独立保证优化或测试表现。
\ans{等价 OR 分数为 $(s_1+s_2)/2$，双极输出为 $2H(u)-1$。指定的一次更新得到 $(1,1,-1)$；编码本身没有普遍准确率保证。}
\qcheck{输入与输出变换都可逆：$x=(s+1)/2$、$z=(t+1)/2$；逐点回代能检验是否保留边界归属。}
\pitfall{双极误分类时 $t-\hat t=\pm2$。若改用误分类更新 $\eta t\tilde s$，步长常数不同，不能在同一数值学习率下声称更新相同。圆圈内阈值为 $\theta$ 时偏置为 $-\theta$。}
\transfer{对一般 $u=w\T x+b$，输入改为 $s=2x-\mathbf1$ 后有 $v=w/2$、$b_s=b+\frac12\sum_iw_i$。换三输入多数函数时仍先逐坐标代入，再判断阈值。}

\problem{Q17}{ADALINE：连续残差的一次 LMS 更新}{正面第3栏／A01、D02}
\qtext{ADALINE 使用 $a=w\T x+b$，目标 $z=1$，单样本损失 $E=\frac12(z-a)^2$。已知 $x=(1,2)\T$、$w=(0.1,-0.1)\T$、$b=0$，$\eta=0.1$。手推梯度并一次同步更新全部参数，计算新输出与新损失。若最后用 $\operatorname{sign}(a)$ 分类，分类正确后是否应停止所有 LMS 更新？}
\route{圈出“\kw{ADALINE/LMS、linear output、squared error}”，查 A01；D02 区分两种误差。更新依据是目标减\kw{线性输出}，不是目标减硬阈值类别。即使题目最后要求分类，训练仍可最小化连续平方残差。}
\given{$x,w\in\R^2$，$a,z,b,E$ 为标量；$e=z-a$。本题有 $1/2$，所以 $\nabla_w E=-ex$、$\partial E/\partial b=-e$，无额外2因子；线性输出斜率为1。}
\solve
先用旧参数前向：
\[
 a=0.1\times1-0.1\times2=-0.1,\quad e=1-(-0.1)=1.1,
 \quad E=\frac12(1.1)^2=0.605.
\]
用链式法则逐项看，例如 $\partial E/\partial w_2=(a-z)\partial a/\partial w_2=-1.1\times2=-2.2$。因此
\[
 \nabla_w E=(-1.1,-2.2)\T,\qquad \frac{\partial E}{\partial b}=-1.1.
\]
同一旧残差驱动所有更新：
\[
 \begin{aligned}
 w_1'&=0.1-0.1(-1.1)=0.21,\\
 w_2'&=-0.1-0.1(-2.2)=0.12,\\
 b'&=0-0.1(-1.1)=0.11.
 \end{aligned}
\]
全部更新完成后再前向：
\[
 a'=0.21\times1+0.12\times2+0.11=0.56,
 \quad e'=0.44,\quad E'=\frac12(0.44)^2=0.0968.
\]
旧分数为负、新分数为正，类别已从−1变成+1，但 $a'=0.56\ne z=1$，平方损失梯度仍非零。分类判对与残差为0是两个条件。
\ans{新参数 $w'=(0.21,0.12)\T,b'=0.11$，新输出0.56，损失从0.605降至0.0968；分类正确并不意味着 LMS 残差已为0。}
\qcheck{直接回代新参数核对输出；也可对 $E$ 作小扰动检查梯度符号。本样本损失下降不代表其他样本的损失也下降。}
\pitfall{若用硬分类差 $1-(-1)=2$ 更新，就改成了另一套规则。若题目给 $E=(z-a)^2$，梯度应多乘2；若 $a=\kappa(w\T x+b)$，梯度还须乘 $\kappa$。}
\transfer{换样本或目标，仍按“线性输出→连续残差→梯度→同步更新→新输出”执行。换成 batch 时，应在同一旧参数处算全体残差，累加梯度后只更新一次。}

\problem{Q18}{Batch、online 与“最小平方仍做不了 XOR”}{正面第3栏／A01；正面第2栏／P02}
\qtext{(a) 只有可训练 bias，模型 $a=b$，初值 $b=0$；两个目标依次为1、2，$\eta=0.1$。分别用 online 单样本半平方、batch 半平方总和、batch 半平方均值完成一轮更新。(b) 对 XOR，增广设计矩阵与目标为
\[
 X=\begin{bmatrix}0&0&1\\0&1&1\\1&0&1\\1&1&1\end{bmatrix},\quad
 z=\begin{bmatrix}-1\\1\\1\\-1\end{bmatrix}.
\]
用 $E=\frac12\norm{z-X\tilde w}^2$ 推导驻点并求解，解释结果是否意味着 XOR 已被分类学会。}
\route{圈出“\kw{online/batch、sum/mean、least squares、stationary}”，查 A01。第一问比较更新时机及缩放；第二问用线性输出平方目标，不能使用感知机误分类更新。P02 用来检查结构能否实现标签。}
\given{(a) 单样本梯度 $\partial E_p/\partial b=b-z_p$。(b) $X\in\R^{4\times3}$，$\tilde w=(w_1,w_2,b)\T\in\R^3$，$z\in\R^4$；$X\tilde w$ 是四个连续预测。}
\solve
(a) Online 更新两次，第二个残差使用更新后的 bias：
\[
 b_1=0+0.1(1-0)=0.1,\qquad b_2=0.1+0.1(2-0.1)=0.29.
\]
Batch 总损失在旧 $b=0$ 处累加，得 $b'=0+0.1[(1-0)+(2-0)]=0.3$。Batch 均值梯度再除样本数2，得 $b'=0.15$。0.29与0.3来自\kw{残差计算时机}不同；0.3与0.15来自\kw{目标缩放}不同。

(b) 先对参数求导：
\[
 \nabla E=X\T(X\tilde w-z)=0
 \quad\Longrightarrow\quad X\T X\tilde w=X\T z.
\]
逐项内积得到
\[
 X\T X=\begin{bmatrix}2&1&2\\1&2&2\\2&2&4\end{bmatrix},\qquad
 X\T z=\begin{bmatrix}0\\0\\0\end{bmatrix}.
\]
三条方程为 $2w_1+w_2+2b=0$、$w_1+2w_2+2b=0$、$2w_1+2w_2+4b=0$。第一条减第二条得 $w_1=w_2=s$；第一条给 $3s+2b=0$；第三条给 $4s+4b=0$，即 $b=-s$。代回得 $s=0$，所以 $\tilde w=0$。

矩阵 $X$ 满列秩，Hessian $X\T X$ 正定，因此这是唯一全局最小值。预测全为0，残差仍为 $(-1,1,1,-1)\T$，$E=2$；无论把零分数统一归到哪一类，都不可能正确分开 XOR。
\ans{(a) Online 得0.29，batch 总和得0.3，batch 均值得0.15。(b) 最小平方唯一最优为 $w_1=w_2=b=0$；优化已解决该目标，结构仍无法表示 XOR 分类。}
\qcheck{把零向量代入正规方程满足驻点；$X\T z=0$ 表明最优残差与设计矩阵列正交。它不要求每个残差都为0。}
\pitfall{不能写“ADALINE只能batch”：课件展示batch版本，online LMS同样可由单样本平方损失求导。驻点一般不自动是最小值；本题依赖凸平方目标和满列秩作结论。}
\transfer{换数据后重新计算 $X\T X$ 与 $X\T z$；若列不独立，最优系数可能不唯一。换总和/均值时相应调整梯度尺度；固定步长online也不保证每轮总损失单调下降。}

\problem{Q19}{Perceptron V.S. ADALINE：可直接写的比较}{正面第3栏／D02、A01；正面第2栏／P01}
\qtext{用完整句子比较 Perceptron 与 ADALINE，必须涉及共同点、输出/损失/更新依据、训练方式、适用任务和限制。再判断：某题给目标 $z=1$、当前线性分数 $u=0.2$；是否“两个模型已预测正确，所以都不再更新”？假设感知机用 $H(u)$，ADALINE 用半平方损失。}
\route{圈出“\kw{compare/difference/common、hard limiter、linear residual}”，查 D02。简答题先选明确比较维度，再组织句子。看到“学习率/更新权重”不足以识别模型，须继续检查输出和损失。}
\given{两者共同有 $u=w\T x+b$；感知机输出 $\hat z=H(u)$，ADALINE 的训练输出为 $a=u$。目标使用0/1编码，$H(0)=1$；$\eta>0$，若 $x\ne0$，非零残差通常产生非零权重更新。}
\solve
可写成以下完整比较：

\kw{共同点：}两者都是用监督目标调整权重的单层前馈模型，均先形成仿射加权和并可含 bias；若最终用一个阈值分类，原输入中的决策边界都是超平面。

\kw{输出与学习依据：}感知机把加权和送入硬阈值，按目标类别与预测类别的差纠错，已分类正确的样本不触发该规则更新。ADALINE 在训练时保留线性连续输出，最小化平方残差；即使类别正确，只要分数尚未达到目标，梯度仍可非零。

\kw{优化与训练方式：}课件所列感知机规则不靠硬阈值输出误差的普通平滑求导；ADALINE/LMS 由可微平方目标的梯度下降得到。两者均可安排逐样本或整批更新；课件“ADALINE only batch”不能作为一般限制。总和与均值损失需协调学习率。

\kw{任务与限制：}感知机用于线性阈值分类，有限可分训练集有标准收敛保证；ADALINE 可用于线性回归，也可对连续分数阈值化作分类。二者在原始特征上的单一仿射分数都无法正确表达 XOR；ADALINE 的最优平方误差不等于零分类错误。

对指定样本，感知机预测 $H(0.2)=1$，类别差0，故无更新。ADALINE 残差 $e=1-0.2=0.8$，因此 $\Delta w=0.8\eta x$、$\Delta b=0.8\eta$，仍更新。
\ans{共同点是监督、单层前馈和仿射加权和；主要差异是硬类别纠错与线性平方残差梯度。当前感知机不更新，ADALINE 仍按残差0.8更新。自查关键词：\kw{hard limiter / linear output / class error / squared residual / supervised / linear boundary}。}
\qcheck{把两套更新中的“误差”明确写出来：$1-H(0.2)=0$ 与 $1-0.2=0.8$，差异立即可见。}
\pitfall{只说“都是神经网络”不足以比较；说“ADALINE用梯度所以能解所有非线性分类”也错误。不要把 batch 误称为“同时把每个参数逐个改一次后重算残差”。}
\transfer{如果题目换成±1标签，重新注明阈值和误差幅度；比较维度不变。若模型先换了非线性特征，要说明线性边界是在特征空间，而不再断言原输入空间只有直线。}

\problem{Q20}{MLP 前向与 XOR 的两种构造}{正面第3栏／B01；正面第2栏／P02}
\qtext{所有阈值单元使用 $H(u)=1$ 当 $u\ge0$，否则0。网络 A 定义
\[
 h_1=H(x_1-x_2-1/2),\quad h_2=H(-x_1+x_2-1/2),\quad
 y=H(h_1+h_2-1/2).
\]
对四个二值输入计算隐藏状态与输出，说明它为何能实现 XOR。网络 B 只用一个隐藏单元 $h=H(x_1+x_2-3/2)$，输出允许直接连接原输入：$y=H(x_1+x_2-2h-1/2)$。检查 B 的四点输出，并回答：若所有隐藏激活改成线性函数，只加更多层能否解决 XOR？}
\route{圈出“\kw{MLP、hidden layer、forward pass、XOR、skip connection}”，查 B01。这里只算固定参数的推理，无需 BP；BP 是求训练梯度的方法。P02 说明单一超平面为什么不够。}
\given{$x=(x_1,x_2)\T$ 是输入；A 有2个非线性隐藏单元和1个输出单元，两层可训练计算。若计入只传值的输入层，画图会有3层，层数必须说明计法。式中 $-1/2$ 是 bias，不是正阈值。}
\solve
网络 A 按“两个隐藏加权和→两个阈值输出→输出层”计算：
\[
\begin{array}{c|rr|cc|r|c}
x&u_{h_1}&u_{h_2}&h_1&h_2&u_y&y\\\hline
(0,0)&-1/2&-1/2&0&0&-1/2&0\\
(0,1)&-3/2&1/2&0&1&1/2&1\\
(1,0)&1/2&-3/2&1&0&1/2&1\\
(1,1)&-1/2&-1/2&0&0&-1/2&0
\end{array}
\]
$h_1$ 检测“$x_1$ 比 $x_2$ 大”，$h_2$ 检测反方向，输出对两种“不同”作 OR。原空间需要两个区域；隐藏空间中正类为01/10，负类为00，用 $h_1+h_2\ge1/2$ 一次分开。

网络 B 的隐藏单元是 AND，四点 $h=0,0,0,1$。输出分数依次为 $-1/2,1/2,1/2,-1/2$，仍得 XOR。它只有两个计算单元，但有\kw{输入直接到输出的跨层连接}，不能把它误画成普通无跨层连接的一个隐藏节点网络。

若隐藏层仅作仿射变换，$h=Ax+c$，输出激活前分数 $v\T h+d=v\T Ax+v\T c+d$ 仍是仿射函数。任意多层仿射变换复合仍仿射；末尾一个阈值仍只有一个超平面，不能突破 XOR 限制。
\ans{A 与 B 都输出 $0,1,1,0$。A 通过非线性隐藏表示把 XOR 变为可分；B 还依赖跨层连接。增加纯线性/仿射层不能解决 XOR。}
\qcheck{四点真值表是本题完整检验；不能只算输入 $(1,0)$ 就宣称网络实现 XOR。A 的隐藏输出维数2与输入维数2相同，非线性改变表示不一定升维。}
\pitfall{普通 BP 不能对硬阈值按 sigmoid 导数求导；本题是固定逻辑网络的前向计算。隐藏层没有题目直接给出的目标，不表示它的活动不能被计算或观察。}
\transfer{换权重或阈值，前向顺序不变，但需重新逐点检查。换为 sigmoid 网络时，输出通常为连续近似值，应分清实际激活与阈值化后的类别，不能把0.9直接写成精确布尔1。}

\problem{Q21}{完整一次 BP：四个参数全部同步更新}{正面第4栏／B02；激活见正面第3栏／B01}
\qtext{一个输入、一个隐藏节点和一个输出节点：
\[
 h=\sig(q),\quad q=wx+b_h,\qquad
 o=\sig(r),\quad r=vh+b_o,\qquad \sig(u)=\frac1{1+e^{-u}}.
\]
取 $x=z=1$，初值 $(w,b_h,v,b_o)=(0,0,2,-1)$，$E=\frac12(z-o)^2$，$\eta=1$，无 momentum。按课件约定 $\delta=-\partial E/\partial(\text{激活前输入})$，完成前向、反向、四个梯度及一次更新。}
\route{圈出“\kw{MLP、sigmoid、half squared error、backpropagation、one update}”，查 B02。B01 提供局部导数。与 LMS 不同，这里需把输出残差经输出激活、连接权及隐藏激活依次传回；直接把残差乘输入会漏因子。}
\given{$x,z,w,b_h,v,b_o,h,o$ 都是标量。$w$ 是输入到隐藏权，$v$ 是隐藏到输出权。本题统一负导数 $\delta_o=-\partial E/\partial r$、$\delta_h=-\partial E/\partial q$，所以参数更新为加 $\eta\delta\times$前层活动。}
\solve
\kw{1. 旧参数前向并记录活动。}
\[
 q=0\times1+0=0,\quad h=\frac12,\qquad
 r=2\times\frac12-1=0,\quad o=\frac12,\quad E=\frac18.
\]
\kw{2. 输出层链式求导。}三个因子为残差方向、局部激活导数、前层活动：
\[
 \delta_o=(z-o)o(1-o)=\frac12\cdot\frac14=\frac18,
\]
\[
 \frac{\partial E}{\partial v}=-\delta_o h=-\frac1{16},\qquad
 \frac{\partial E}{\partial b_o}=-\delta_o=-\frac18.
\]
\kw{3. 隐藏层接收下游贡献。}此时仍用旧连接权 $v=2$：
\[
 \delta_h=h(1-h)v\delta_o=\frac14\cdot2\cdot\frac18=\frac1{16},
\]
\[
 \frac{\partial E}{\partial w}=-\delta_h x=-\frac1{16},\qquad
 \frac{\partial E}{\partial b_h}=-\delta_h=-\frac1{16}.
\]
\kw{4. 梯度全部算完后统一更新。}
\[
 \boxed{(w',b_h',v',b_o')=\left(\frac1{16},\frac1{16},\frac{33}{16},-\frac78\right).}
\]
例如 $v'=2-(-1/16)=33/16$。新参数用于下一次前向，而不用于重写本次隐藏梯度。
\ans{旧 $h=o=1/2$、$E=1/8$；$\delta_o=1/8,\delta_h=1/16$。四个梯度依 $(w,b_h,v,b_o)$ 顺序为 $(-1/16,-1/16,-1/16,-1/8)$，更新参数如框中所示。}
\qcheck{新 $q'=1/8$，故 $h'=\sig(1/8)$；$o'=\sig((33/16)\sig(1/8)-7/8)\approx0.5549$，$E'\approx0.0990<0.125$。小扰动中心差分也可核对四个旧梯度。}
\pitfall{若先更新 $v$ 再算 $\delta_h$，就混用两套参数。若把 $\delta$ 定义为正导数，公式符号会整体变化；不能正导数配“加 $\eta\delta$”。bias 的梯度无需再乘 $h$ 或 $x$。}
\transfer{多输出时，隐藏节点把所有下游 $w\delta$ 相加；batch 时先在共同旧参数处累加各样本梯度。换激活应重算局部导数；换成不含 $1/2$ 的平方损失，全部梯度多乘2。}

\problem{Q22}{隐藏层的分支求和与 momentum}{正面第4栏／B02、B03}
\qtext{某 logistic 隐藏节点当前活动 $h=0.8$，有两个下游连接，旧权分别 $v_1=0.5,v_2=-0.3$，下游负导数 deltas 为 $\delta_1=0.1,\delta_2=-0.2$。该节点一个输入为 $x=-2$，对应旧权 $w=0.4$。先求隐藏 $\delta_h$ 与 $\partial E/\partial w$。再用 $\eta=0.05$、动量系数 $\mu=0.8$、上一步权重增量 $\Delta w_{\rm prev}=0.01$，求新增量与新权重。解释动量是否保证找到全局最小。}
\route{圈出“\kw{hidden、two downstream paths、momentum、previous increment}”，查 B02/B03。隐藏节点没有直接目标，不能自行编一个 $z-h$；必须汇总所有下游链式贡献。动量还需读取上一步\kw{增量}。}
\given{$\delta_j=-\partial E/\partial u_j$；$\delta_h=h(1-h)\sum_jv_j\delta_j$，$\partial E/\partial w=-\delta_h x$。动量规则为 $\Delta w=-\eta\partial E/\partial w+\mu\Delta w_{\rm prev}$，$w'=w+\Delta w$。}
\solve
先算两条分支的加权和：
\[
 v_1\delta_1+v_2\delta_2=0.5(0.1)+(-0.3)(-0.2)=0.11.
\]
隐藏局部导数为 $h(1-h)=0.8\times0.2=0.16$，所以
\[
 \delta_h=0.16\times0.11=0.0176,\qquad
 \frac{\partial E}{\partial w}=-(0.0176)(-2)=0.0352.
\]
虽然隐藏 delta 为正，这条权重的\kw{负输入}使梯度为正；无动量梯度下降会减小该权重。
\[
 \Delta w=-0.05(0.0352)+0.8(0.01)=-0.00176+0.008=0.00624,
 \qquad w'=0.40624.
\]
本次增量反而为正，因为前次方向的累积项更大。动量可平滑方向变化、减轻狭长误差谷中的振荡，但不是全局搜索证明；它也不保证每一步目标下降。
\ans{$\delta_h=0.0176$，梯度0.0352，新增量0.00624，新权重0.40624。动量保留历史运动方向，不保证全局最优。}
\qcheck{删去动量即得到 $w'=0.39824$；与原结果比较可检查两项符号。隐藏分支若只保留第一项，会误得 $\delta_h=0.008$，明显漏掉第二条有效路径。}
\pitfall{不能用 $\mu w$ 代替 $\mu\Delta w_{\rm prev}$。同一步的分支权应为旧参数；负连接权与负 delta 的乘积为正，不可凭箭头符号直觉跳过乘法。}
\transfer{增加下游节点只需扩大求和；换为 tanh 隐藏节点须将局部导数改成 $1-h^2$。若题目给的“动量变量”是速度而非参数增量，先把其递推和符号与本规则对应，避免混用定义。}

\problem{Q23}{激活函数纠错：少一个尺度就改了梯度}{正面第3栏／B01；反传见正面第4栏／B02}
\qtext{判断并改正 $g(u)=(1-e^{-u})/(1+e^{-u})=\tanh(u)$。推导 $g'(u)$ 并用输出 $g(u)=0.6$ 计算局部导数。另有斜率为2的 logistic $a(u)=1/(1+e^{-2u})$，求 $a'(0)$。最后解释：为何硬阈值不能直接沿用这些导数，sigmoid 的饱和会带来什么训练问题？}
\route{圈出“\kw{sigmoid/tanh、derivative、slope、correct the formula}”，查 B01。先看\kw{完整指数和尺度}，不能因为两条曲线都为S形就套同一个导数；B02 中该局部导数会乘入整条链式梯度。}
\given{$\sig(u)=(1+e^{-u})^{-1}$，$\sig'(u)=\sig(u)(1-\sig(u))$；$\tanh(u)=(1-e^{-2u})/(1+e^{-2u})$。这里 $u$ 为激活前输入，$g,a$ 为活动值，均为标量。}
\solve
将 $u/2$ 代入标准 tanh 定义：
\[
 g(u)=\tanh(u/2)=2\sig(u)-1.
\]
由链式法则或商法则可得
\[
 g'(u)=\frac12[1-\tanh^2(u/2)]
 =\frac12(1-g(u)^2).
\]
所以当 $g=0.6$ 时，$g'=\frac12(1-0.36)=0.32$。若误用标准 $\tanh(u)$ 导数，会得到0.64，为正确值的两倍。

斜率2的 logistic 是 $a(u)=\sig(2u)$，故
\[
 a'(u)=2a(u)(1-a(u)),\qquad a(0)=\frac12,\qquad a'(0)=\frac12.
\]
普通 $\sig(u)$ 在0处斜率为1/4，标准 $\tanh(u)$ 在0处为1，$\tanh(u/2)$ 在0处为1/2，可作为三个公式的快速尺度检查。

硬阈值在非边界处导数为0、跳变处不可按普通微分求导，不能直接给它套 $a(1-a)$。Logistic/tanh 在两端饱和，导数接近0，链式传播时小导数连乘可使早层梯度很小，导致学习缓慢。
\ans{原等式应改为 $g(u)=\tanh(u/2)$，导数 $\frac12(1-g^2)$；指定活动处导数0.32；斜率2的 logistic 在0处导数0.5。}
\qcheck{在 $u=0$ 附近取很小的正负扰动，用函数值差商检查斜率；标准 tanh 与本题 $g$ 输出零点相同，但斜率不同。}
\pitfall{饱和区导数\kw{接近}0，不表示 sigmoid 处处导数为0。分段线性压缩函数在折点处也需另定导数，不能把“连续”直接当“处处可微”。}
\transfer{一般 $\sig(cu)$ 多乘 $c$；一般 $\tanh(cu)$ 的导数为 $c(1-a^2)$。换激活只改变局部导数因子，前向记录、反向链式和同步更新的框架保持不变。}

\problem{Q24}{表示能力、优化成功与泛化是三件事}{正面第4栏／B03；模型结构见正面第3栏／B01}
\qtext{有人说：“MLP 是通用逼近器，Kolmogorov 也说明多变量函数可分解，所以固定一个很小的 sigmoid 网络、持续运行 BP，一定会找到全局最优并对新数据正确。”逐项回应。再看真实规律 $f(x)=x$，训练数据只有 $(0,0),(1,1)$；候选为 $g_1(x)=x$ 与 $g_2(x)=x+4x(1-x)$。比较训练误差及 $x=1/2$ 处预测。若训练损失持续下降而验证损失先降后升，应怎样使用验证集与测试集？}
\route{圈出“\kw{universal approximation、exist/find/generalize、overfitting、validation}”，查 B03。这里不是再算一次梯度，而是检查结论声称解决了哪一个问题：存在参数、算法找到参数、未见数据有效。}
\given{表示能力讨论模型家族及容量；优化讨论选定目标、算法和初值；泛化讨论训练与模型选择之外的数据。算术小问统一用单点半平方损失 $\frac12(f-g)^2$，输入域为 $[0,1]$。}
\solve
\begin{enumerate}
\item \kw{表示：}通用逼近结论在适当激活、容量和紧域上的连续目标等条件下，允许选择足够的隐藏单元，使近似误差任意小。它不是固定某个很小网络对任意函数精确相等的结论。Kolmogorov 的存在性分解允许合适的一元函数，不能未经证明把它们全部换成一个指定的小 sigmoid 网络。
\item \kw{优化：}参数存在不代表 BP 从任何初值都能找到。一般多层训练目标非凸，可能存在局部极小、鞍点或平坦区域；驻点也不自动是全局最优。Momentum 能改善某些轨迹，但没有解除上述限制。
\item \kw{泛化：}有限训练点不能唯一确定点间行为。低训练误差没有独自证明对新样本的表现，通用逼近存在性也没有提供这种保证。
\end{enumerate}
本例中，两候选在0、1处都与真实目标相同，训练损失都为0。但在 $x=1/2$，
\[
 f=\frac12,\quad g_1=\frac12,\quad g_2=\frac12+4\cdot\frac12\cdot\frac12=\frac32.
\]
故该点 $g_1$ 损失0，$g_2$ 损失 $\frac12(1)^2=1/2$。这是零训练误差仍可泛化较差的明确例子。

训练/验证损失分离提示过拟合或数据划分问题。用验证集选择训练轮次和模型设置，保留验证表现较好的参数；可比较容量、正则化等设置。测试集应保留为最终评价，不反复据测试结果调参后仍称其为独立检验。
\ans{存在性、优化、泛化分别需要不同证据，三者不能互相替代。本例两模型训练误差均为0，但在 $1/2$ 处 $g_2$ 预测1.5、半平方损失0.5；继续降低训练误差不能保证恢复泛化。关键词：\kw{capacity / nonconvex / validation / overfitting}。}
\qcheck{先用训练点验证“两者都拟合”，再用未参与拟合的点验证差异；不要只凭曲线弯曲与否定义过拟合。}
\pitfall{“网络大”并非过拟合的充分定义；“总损失驻点”也不是“每个样本残差为0”。不能从盲人摸象插画推出课件没有给出的数学定理。}
\transfer{若改成不连续目标或换域，须重新检查逼近定理条件。模型选择程序仍应区分训练、验证、测试；测试好也只支持已进行的评价范围，不能保证任意未来分布。}

\problem{Q25}{RBF 特征值、中心和宽度}{正面第5栏／R01}
\qtext{Gaussian RBF 定义为 $\phi_j(x)=\exp[-\norm{x-c_j}^2/\sigma_j^2]$。取 $c_1=(0,0)\T,c_2=(1,1)\T$，两个宽度都为1。计算 $x=(1,0)\T$ 与 $x=(0,0)\T$ 的特征向量。若同样距离 $r=1$ 的一个基函数宽度由1改为2，响应如何改变？若题目改用 $\exp[-r^2/(2\sigma^2)]$，$r=\sigma=1$ 时又是多少？说明“径向”与“Gaussian 局部响应”的区别。}
\route{圈出“\kw{RBF、center、distance、width、Gaussian}”，查 R01。先核对分母约定，再算平方距离和指数。这里是在计算显式特征；仅见 Gaussian 并不能认定题目正在训练 SVM。}
\given{$x,c_j\in\R^2$，$r_j=\norm{x-c_j}\ge0$，$\sigma_j>0$；$\phi(x)=(\phi_1,\phi_2)\T\in\R^2$。输出权重尚未给出，因此本题只算特征，不能自行把特征和称为预测。}
\solve
对 $x=(1,0)\T$，两个平方距离均为1：
\[
 \norm{x-c_1}^2=1^2+0^2=1,\qquad
 \norm{x-c_2}^2=0^2+(-1)^2=1,
 \quad\phi(x)=(e^{-1},e^{-1})\T\approx(0.3679,0.3679)\T.
\]
对 $x=(0,0)\T$，到两个中心的平方距离为0、2，所以
\[
 \phi(x)=(1,e^{-2})\T\approx(1,0.1353)\T.
\]
固定 $r=1$ 时，宽度由1变2使响应由 $e^{-1}\approx0.3679$ 变为 $e^{-1/4}\approx0.7788$。宽度越大，离中心同样远的位置仍可获得较大响应；不是峰值变高，中心处峰值仍为1。

若分母给 $2\sigma^2$，则响应为 $e^{-1/2}\approx0.6065$，应照题使用；两个宽度参数的数值含义不同。

“径向”表示只依赖距离 $r$；同中心、同半径的位置有相同响应。课件还列 $r^p$（$p=1,3,\ldots$）、薄板 $r^2\log r$、多二次 $\sqrt{r^2+\sigma^2}$ 及其倒数。它们也径向，但不都像 Gaussian 那样远离中心趋于0；薄板式在 $r=0$ 按极限取0。
\ans{两个输入的特征分别为 $(e^{-1},e^{-1})\T$ 与 $(1,e^{-2})\T$。宽度1变2，指定响应由0.3679变0.7788；分母 $2\sigma^2$ 的响应为0.6065。Gaussian 的局部钟形是特定基函数的性质。}
\qcheck{中心处平方距离0、响应1；有限距离时 Gaussian 响应应在 $(0,1]$。两个中心等距时两响应相等。}
\pitfall{题目已经给 $\norm{x-c}^2$，不要再平方成距离四次方；$r$ 是距离，$r^2$ 才是平方距离。宽度参数要正；不要把分母里的 $\sigma^2$ 与 $2\sigma^2$ 混用。}
\transfer{换中心会改变每点到中心的距离；换宽度只改指数尺度。中心与宽度固定后可查 R01 解线性输出；若同时优化中心/宽度，整体训练就不再只是线性最小二乘。}

\problem{Q26}{固定 RBF 特征后求输出权重}{正面第5栏／R01；矩阵消元见正面第1栏／M03}
\qtext{一维 RBF 网络有中心 $c_1=0,c_2=1$，$\sigma^2=1/\log2$，$\phi_j(x)=e^{-(x-c_j)^2/\sigma^2}$。输出 $f(x)=a_1\phi_1(x)+a_2\phi_2(x)$，\kw{不另加 bias}。要求在 $(x,z)=(0,1),(1,0)$ 两点精确插值。构建特征矩阵、消元求 $a_1,a_2$，预测 $x=1/2$。解释为何若额外加入自由 bias，不能继续声称两条方程唯一决定三个未知数。}
\route{圈出“\kw{fixed centers/width、linear output、interpolation、solve weights}”，查 R01。本题未知量是输出系数，特征已固定，适合线性方程求解；无需对中心做 BP，也不需要 SVM 的间隔约束。}
\given{两样本 $N=2$、两基函数 $h=2$，$\Phi\in\R^{2\times2}$，$a=(a_1,a_2)\T\in\R^2$，$z=(1,0)\T$。矩阵第 $i$ 行是第 $i$ 个样本的特征，求解 $\Phi a=z$。}
\solve
距离平方为0时响应1，为1时响应 $e^{-\log2}=1/2$，故
\[
 \Phi=\begin{bmatrix}1&1/2\\1/2&1\end{bmatrix},\qquad
 \begin{bmatrix}1&1/2&\vline&1\\1/2&1&\vline&0\end{bmatrix}.
\]
作行变换 $R_2\leftarrow R_2-\frac12R_1$：
\[
 \begin{bmatrix}1&1/2&\vline&1\\0&3/4&\vline&-1/2\end{bmatrix}.
\]
回代得 $a_2=-2/3$，再由第一行 $a_1=1-\frac12(-2/3)=4/3$。

预测点到两中心的平方距离均为 $1/4$，所以
\[
 \phi_1(1/2)=\phi_2(1/2)=e^{-(\log2)/4}=2^{-1/4},\qquad
 f(1/2)=\left(\frac43-\frac23\right)2^{-1/4}\approx0.5606.
\]
虽然每个 Gaussian 响应非负，输出权重可以为负。

若另加自由 bias $b$，矩阵要增广一列1，未知变为 $(a_1,a_2,b)$。两行三列不可能满列秩；存在插值解时通常不唯一，需要更多独立数据或另外给定选择/正则化原则。不能把原二阶矩阵的解直接当新模型唯一解。
\ans{$a=(4/3,-2/3)\T$，$f(1/2)=\frac23\,2^{-1/4}\approx0.5606$。固定特征后，对输出权重是线性求解，对原输入仍是非线性函数。}
\qcheck{$f(0)=4/3-(2/3)(1/2)=1$，$f(1)=(4/3)(1/2)-2/3=0$；$\det\Phi=3/4\ne0$，本题插值系数唯一。}
\pitfall{矩阵行是样本、列是特征；中心数量与样本数量概念不同。只因本题方阵可逆，不能推出所有 RBF 训练都能直接求逆；一般近似问题是 $\min_a\norm{z-\Phi a}^2$。}
\transfer{换目标仅改右端；换中心/宽度须重算特征矩阵。若数据更多，通常解最小二乘而非强制插值；若特征列相关，正规方程可能非唯一，应按题目指定的伪逆或正则化规则处理。}

\problem{Q27}{Cover 思想：XOR 换表示后怎样可分？}{正面第5栏／R01；线性不可分见正面第2栏／P02}
\qtext{XOR 标签为00/11负类，01/10正类。按新课件中心顺序固定 $c_1=(1,1)\T,c_2=(0,0)\T$，$\phi_j(x)=e^{-\norm{x-c_j}^2}$。重算四点特征；检验 $g(x)=0.9-\phi_1(x)-\phi_2(x)$ 能否分开两类。写出带 bias 的 $\phi$ 可分条件，并纠正“任何升维映射都会可分、会提高泛化”。}
\route{圈出“\kw{Cover、$\phi$-separable、XOR、feature space}”，查 R01。先计算映射再检查线性分数；原空间 XOR 不可分不妨碍在合适新特征中可分。此题中心顺序与 Q25 不同，不能直接抄上一题列顺序。}
\given{$x,c_j\in\R^2$，$\phi(x)\in\R^2$；输出分数的权重为 $a=(-1,-1)\T$，bias为0.9。需要齐次形式时用 $\tilde\phi=(\phi_1,\phi_2,1)\T$、$\tilde a=(-1,-1,0.9)\T$。}
\solve
令 $q=e^{-1}\approx0.367879$、$r=e^{-2}\approx0.135335$。先按两个中心顺序计算平方距离：
\[
\begin{array}{c|cc|cc|c}
x&\norm{x-c_1}^2&\norm{x-c_2}^2&\phi_1&\phi_2&g\\\hline
(0,0)&2&0&r&1&0.9-(1+r)\\
(0,1)&1&1&q&q&0.9-2q\\
(1,0)&1&1&q&q&0.9-2q\\
(1,1)&0&2&1&r&0.9-(1+r)
\end{array}
\]
负类分数约为 $-0.235335$，正类分数约为 $0.164241$，均严格符合符号。因此在新坐标中，直线 $\phi_1+\phi_2=0.9$ 分开了两类。

本例两个不同的正类输入01与10映到同一个特征点，但标签相同，并不妨碍分类。新映射还是二维到二维；起作用的是\kw{非线性改变表示}，不只是增加坐标数量。原输入中的边界变成两个指数函数之和等于0.9，一般不再是一条直线。

有 bias 的 $\phi$ 可分性指存在 $a,b$ 使正类 $a\T\phi(x)+b>0$、负类 $a\T\phi(x)+b<0$；用增广常数坐标可写成齐次形式。课件以 Cover 为标题主要给出这个定义及表示思想，不应声称已展示了完整定理证明。

反例：把任意输入映为零向量，再增加任意多个零坐标，异类样本仍完全重合，任何确定的线性输出都无法区分。即使训练集变可分，也没有因此自动获得测试集保证。
\ans{指定映射后 XOR 严格可分；00/11分数约−0.2353，01/10约0.1642。适当非线性表示可以提高可分性，任意映射或任意升维没有此保证，泛化还需独立评价。}
\qcheck{与课件p146一致时，00必须映为 $(e^{-2},1)$；p147原表及图交换了两特征列。中心、列、图轴与权重必须同步，不能拼接不一致的定义。}
\pitfall{线性可分须说清在哪个空间。不要把本题0.9分类线称为已求得 SVM 最大间隔线；本题只验证符号，不做间隔优化。}
\transfer{交换两中心时，应同时交换特征列和对应输出权重；本题权重相同，预测恰好不变。若权重不同或宽度改变，须重新计算分数。分类正确也不等于精确插值到0/1目标。}

\problem{Q28}{FLN 与 ELM：特征从哪里来，求哪些参数？}{正面第5栏／R02；比较见 D04}
\qtext{(a) 二值输入的 FLN 使用固定特征 $(x_1,x_2,x_1x_2)$，证明输出 $f=x_1+x_2-2x_1x_2$ 精确等于 XOR，说明训练的是哪些量。(b) 基本 ELM 的一次初始化已给定并固定隐藏特征 $\phi_1(x)=\sig((\log3)x)$、$\phi_2(x)=\sig(- (\log3)x)$，$\sig(u)=(1+e^{-u})^{-1}$。不加输出 bias，输出 $f=\beta_1\phi_1+\beta_2\phi_2$；样本为 $(x,z)=(0,0),(1,1)$。构建 $H\beta=z$，求输出权并预测 $x=-1$。解释为什么 ELM 不是“没有隐藏层”。}
\route{圈出“\kw{functional link、fixed expansion、random fixed hidden、output weights}”，查 R02。FLN 的函数项预先指定；基本 ELM 隐藏权和 bias 初始化后固定。此时只有输出系数需用标签拟合，不能把隐藏参数也当本次线性方程的未知数。}
\given{(a) 原输入 $x\in\R^2$，展开特征 $\phi\in\R^3$，系数为 $(1,1,-2)$，bias为0。(b) $N=h=2$，$H_{ij}=\phi_j(x_i)$，$H\in\R^{2\times2}$，$\beta\in\R^2$，$z=(0,1)\T$；隐藏斜率 $\pm\log3$ 与隐藏 bias 0 已固定。}
\solve
(a) 逐点代入：$f(0,0)=0$、$f(0,1)=1$、$f(1,0)=1$、$f(1,1)=2-2=0$。乘积特征记录两输入共同出现的作用。固定展开后，输出对\kw{系数}线性，对\kw{原输入}可以非线性；学习时拟合输出系数及题目允许的 bias，不必把函数项本身也当可训练隐藏层。

(b) $\sig(0)=1/2$，$\sig(\log3)=3/4$、$\sig(-\log3)=1/4$，所以
\[
 H=\begin{bmatrix}1/2&1/2\\3/4&1/4\end{bmatrix},\qquad
 H\beta=\begin{bmatrix}0\\1\end{bmatrix}.
\]
将两式分别乘2、4，得 $\beta_1+\beta_2=0$、$3\beta_1+\beta_2=4$；相减得 $2\beta_1=4$，所以 $\beta_1=2,\beta_2=-2$。在 $x=-1$，特征为 $(1/4,3/4)$，故 $f(-1)=2/4-6/4=-1$。

基本 ELM 含一个隐藏层，通常随机产生并固定隐藏输入权和 bias，再用最小二乘/伪逆等方法求输出权。课件“one-layer”不能解读成没有隐藏节点；\kw{一层可拟合权重}与\kw{一个隐藏层}不是相同的计层口径。固定 $H$ 的线性平方子问题没有坏的局部极小，但特征质量、秩、节点数和泛化仍需判断。
\ans{FLN 的固定乘积特征使 XOR 对系数线性可表示；ELM 本题 $H$ 如上，$\beta=(2,-2)\T$，$f(-1)=-1$。两者主要区别在特征来源，不在“一个非线性、另一个线性”的二分。}
\qcheck{$H(2,-2)\T=(0,1)\T$；$\det H=-1/4\ne0$，本题输出系数唯一。预测−1只是该模型在给定初始化和数据下的结果，不能把它当已获保证的真实规律。}
\pitfall{“随机”不表示推理时重新抽隐藏权。若隐藏权在输出训练期间也持续优化，固定特征线性求解的前提已改变；若 $H$ 非方阵或非满秩，不可无条件写 $H^{-1}$。}
\transfer{换标签仅改右端；换初始化、激活或节点数须重建 $H$。一般可写 $\hat\beta=H^+z$，含额外输出 bias 时须在 $H$ 加一列1并重新检查秩；该线性求解不自动替你选择好的特征。}

\problem{Q29}{MLP V.S. RBF：模型与训练方法分开答}{正面第5栏／D03、R01；正面第4栏／B02}
\qtext{比较 MLP 与 RBF 网络的共同点、隐藏特征、输出、训练方式、任务与限制。有人写“BP、MLP、RBF 是三种并列网络；MLP 一定慢而且不能回归，RBF 一定只做无监督聚类”，请纠正。某任务已给定合适中心和宽度，需要快速拟合连续输出，应怎样说明选择理由及保留条件？}
\route{圈出“\kw{MLP versus RBF、common/difference、choose model、BP}”，查 D03。按架构和参数训练角色比较；不要只列“都是AI”，也不要把求导方法与模型名称排成同类。}
\given{典型 MLP 隐藏活动为 $a_j=s(v_j\T x+b_j)$；Gaussian RBF 隐藏活动为 $\phi_j=e^{-\norm{x-c_j}^2/\sigma_j^2}$。比较的“局部性”针对这里的 Gaussian RBF，不是所有径向基都衰减。输出形式和损失应按任务明确。}
\solve
\kw{共同点：}两者都是前馈参数模型，都能先产生非线性中间特征，再形成预测；适当设计输出和损失后，均可用于分类或回归。模型具有表示能力，不等于已优化成功或已泛化。

\kw{隐藏表示差异：}典型 MLP 节点先作仿射投影再激活，方向、偏置和多层组合可由训练学习；常见 sigmoid 单元对整个投影方向产生响应。Gaussian RBF 节点按到中心的距离局部响应，中心决定位置，宽度决定作用范围；输出是径向特征的线性组合。

\kw{训练差异：}MLP 常用 BP 计算所有可训练层的梯度，再配合梯度下降等优化。RBF 常先选中心/宽度，再固定特征求线性输出权；中心选择可以采用无监督步骤，但输出通常使用监督目标，中心也可以另外监督优化。固定特征时的输出平方问题是凸的；同时训练中心/宽度时，不能声称全问题仍只需一次线性求解。

\kw{任务与限制：}MLP 适合学习层次或组合表示，但训练非凸，架构、初始化与正则化影响结果；RBF 的局部表示直观，固定特征后输出求解直接，但需要足够且合适的中心与宽度，矩阵条件和样本覆盖影响效果。无法仅凭名称保证谁始终更快、更准。

纠错：BP 是链式计算梯度的算法，不是与 MLP/RBF 同层级的网络架构。MLP 可以回归，RBF 网络也并非只有无监督聚类。对所给任务，可选\kw{固定中心和宽度的 RBF+线性输出平方拟合}，理由是隐藏特征已定，输出可由线性最小二乘高效拟合；仍须检查特征覆盖、矩阵秩和独立验证表现。
\ans{两者都可非线性前馈预测；典型 MLP 学习仿射激活的组合，Gaussian RBF 使用到中心的局部距离特征。BP 是梯度方法；RBF 固定特征后可直接拟合线性输出，但不保证所有中心设置都好。关键词：\kw{affine activation / radial distance / local response / fixed features / gradient}。}
\qcheck{比较答案同时出现共同点、同一维度的差异、任务和限制；仅写“一个全局，一个局部”会省掉训练条件与应用。}
\pitfall{前馈指连接方向无环，不表示训练时不能从输出向前层反传导数。误差反传不会把 MLP 变成循环网络。}
\transfer{若换成不同 MLP 输出激活或非 Gaussian 径向函数，须重述相应性质；共同的“特征→输出”结构及表示/优化/泛化区分仍可使用。训练速度应结合数据规模、特征数和求解器讨论。}

\problem{Q30}{RBF V.S. FLN V.S. ELM：共同线性输出的条件}{正面第5栏／D04、R01、R02}
\qtext{比较 RBF 网络、FLN 与基本 ELM 的共同结构、特征来源、固定与训练参数、适用任务及限制。判断：(i) 固定 RBF 中心/宽度后求输出权；(ii) 训练中持续移动 RBF 中心；(iii) 固定 FLN 乘积展开后求输出权；(iv) ELM 每次预测都重新随机隐藏权。哪些保留“固定特征线性输出求解”的前提？解释“快速训练、消除局部极小”能说明什么，不能说明什么。}
\route{圈出“\kw{RBF/FLN/ELM、feature source、which parameters、local minima}”，查 D04。模型比较要追踪特征矩阵是否随未知量改变；只看都含非线性函数不足以判断优化是否为线性平方问题。}
\given{共同写法为 $f(x)=a\T\phi(x)+b$；样本排成 $\Phi_{ij}=\phi_j(x_i)$。当 $\Phi$ 固定、目标为平方残差时，输出参数可作线性最小二乘；若含 bias，增广一列1。}
\solve
\kw{共同点：}三者都可显式计算非线性特征，再对输出系数作线性拟合；按输出和目标设计，可用于分类或回归。对原输入非线性与对输出参数线性可以同时成立。

\kw{RBF：}特征依赖到中心的距离，Gaussian 还需正宽度；中心、宽度、基函数与个数需先选或训练。固定这些特征后拟合输出权和可选 bias；其限制包括中心覆盖、宽度和矩阵条件。

\kw{FLN：}特征是原输入加预设函数项，如乘积、平方；展开函数固定后学习输出系数。它省去另行训练隐藏表示，但函数类型与展开规模需要选择，预设项可能不能充分表达任务规律。

\kw{基本 ELM：}隐藏输入权和 bias 随机初始化后固定，由隐藏响应构成 $H$，主要拟合输出权；隐藏节点数和激活仍要选择。固定的一次随机表示可能质量较差或秩不足，随机并不自动等于可分或泛化好。

逐项判断：(i) 保留；(ii) 不保留，因为移动中心改变每个距离和 $\Phi$；(iii) 保留；(iv) 错误，预测应使用训练时固定的隐藏参数，否则已学习的输出系数对应的特征坐标改变。

固定特征且采用线性输出平方损失时，目标为凸二次函数；其驻点是全局最小，满列秩时输出系数唯一，否则可有多个同等最小值。“无坏局部极小”只适用于这个子问题，不替代特征选择，不保证整个中心/宽度训练凸，也不保证最小训练误差等于最佳泛化。
\ans{三者共同是“显式特征→线性输出”；RBF 按距离，FLN 用预设函数，ELM 用随机固定隐藏响应。(i)(iii) 保留固定特征前提；(ii) 改变特征；(iv) 破坏训练/推理一致性。关键词：\kw{centers/widths / prescribed functions / random fixed hidden / linear least squares}。}
\qcheck{检查未知量是否只出现在最后的 $a,b$；若未知中心、宽度或隐藏权出现在 $\phi$ 内，就不能把整题当固定矩阵求一次线性输出。}
\pitfall{“对参数线性”不等于“对输入线性”。固定特征的平方问题即使凸，也可能不唯一；输出无坏局部极小不等于整个模型没有选择困难。}
\transfer{换数据时保持特征机制定义；换中心、展开项或 ELM 初始化必须重建矩阵并重求输出权。若目标换成其他损失或输出再加非线性，须重新核对是否仍可用线性最小二乘。}

\problem{Q31}{间隔、距离与缩放：哪个量真的改变？}{背面第1栏／S01}
\qtext{一组标签为$y_i\in\{-1,+1\}$的训练样本，给定分离函数
$g(x)=w\T x+b$，其中$w=(3,4)\T,b=-5$；已知$\min_i y_i g(x_i)=1$。求法向量、规范间隔带宽度，以及测试点$x=(1,1)\T,y=+1$的函数间隔和带符号几何间隔。再把$w,b$同时乘2：分类面、该点距离及间隔带宽度如何变化？}
\route{圈出“标签$\pm1$、间隔、距离、缩放”。S01把函数值与真实距离分开；这是分类几何题，尚不需要SVR的连续目标或LS-SVM的平方残差系统。}
\given{$w\in\R^2$，$b\in\R$；$m_i=y_i g_i$是函数间隔，$m_i/\norm w$是带符号几何间隔。题目给出的最小函数间隔1说明当前参数已规范化。}
\solve
先求$\norm w=\sqrt{3^2+4^2}=5$，指向$g>0$一侧的单位法向量为$n=(3/5,4/5)\T$。
规范边界为$g=+1$与$g=-1$，各距决策面$1/5$，全带宽为
\[
M=\frac2{\norm w}=\frac25.
\]
对测试点，先代入函数，再乘标签，最后除范数：
\[
g(1,1)=3+4-5=2,\qquad m=(+1)\times2=2,\qquad
\rho=\frac{m}{\norm w}=\frac25.
\]
它被判为$+1$，位于规范正类间隔边界之外。

令$w'=2w,b'=2b$，则$g'=2g$且$\norm{w'}=10$。决策面$g'=0$与$g=0$相同，正负号相同；测试点的函数间隔变为4，而几何间隔仍为$4/10=2/5$。
原来的两条间隔边界变成$g'=\pm2$，所以同一几何间隔带的宽度为$4/10=2/5$。
新参数的最小训练函数间隔是2，\textbf{不能}把它误当成规范化参数，再直接写$2/\norm{w'}=1/5$。

\ans{单位法向量为$(3/5,4/5)\T$；规范带宽$2/5$；测试点函数间隔2、几何间隔$2/5$。正比例缩放改变函数值，不改变分离面、分类或几何距离。}
\qcheck{点到直线$3x_1+4x_2-5=0$的距离为$|3+4-5|/5=2/5$，与上式一致。}
\pitfall{$w\T x+b$带偏置，严格说是仿射函数。总间隔宽度$2/\norm w$与单侧距离$1/\norm w$不同；前者还要求使用规范边界$g=\pm1$。}
\transfer{更换$w,b,x$时，先重算范数与函数值；若题目给出最小函数间隔$c>0$，几何带宽为$2c/\norm w$，或先除以$c$规范化。}

\problem{Q32}{从两类样本构建并解出硬间隔SVM}{背面第1栏／S01}
\qtext{两个一维训练样本为$(x_1,y_1)=(-1,-1)$、$(x_2,y_2)=(1,+1)$。使用$g(x)=wx+b$、标签$\pm1$和硬间隔规范化。自行写目标和约束，用代数求最优$w,b$，并给出间隔宽度及$x=1/2$的类别。}
\route{“两类标签、硬间隔、求最优分界”指向S01。目标是最大几何间隔，等价于最小化范数；不是把分类标签当连续回归目标做平方误差拟合。}
\given{$x_i,w,b$均为标量，未知量是$w,b$。规范约束为$y_i(wx_i+b)\ge1$。不使用松弛变量，因为本题要求硬间隔且数据可分。}
\solve
最大化$2/|w|$可等价改写为
\[
\min_{w,b}\frac12w^2,
\qquad\text{s.t.}\quad
(-1)(-w+b)\ge1,\quad (+1)(w+b)\ge1.
\]
也就是
\[
w-b\ge1,\qquad w+b\ge1
\quad\Longleftrightarrow\quad w\ge1+|b|.
\]
由此$w\ge1$。目标在可行范围内随$w$增大，因此最小可行值是$w=1$；此时两条约束要求$b\le0$且$b\ge0$，故$b=0$。
\[
g(x)=x,\qquad J_{\min}=\frac12,
\qquad M=\frac2{|1|}=2.
\]
两个样本分别满足$g(-1)=-1,g(1)=1$，均位于间隔边界。
对$x=1/2$，$g=1/2>0$，预测类别为$+1$。测试点可以在间隔带内；训练约束并不要求所有未来测试点都满足$yg\ge1$。

\ans{硬间隔问题为$\min\tfrac12w^2$，约束$w-b\ge1,w+b\ge1$；最优解$(w,b)=(1,0)$，总间隔宽度2，$x=1/2$判为$+1$。}
\qcheck{回代后两点的$y_i g_i$均为1，约束可行；由$w\ge1+|b|$可知不存在范数更小的可行解，因此不是仅找到一个驻点。}
\pitfall{漏掉标签$y_i$会让负类约束写反。最小化$\tfrac12w^2$是优化目标，不能写成“所有SVM量都最大化”。}
\transfer{若改为$x_1=-a,x_2=a$且$a>0$，同样的代数步骤给$w=1/a,b=0$、带宽$2a$。若同一个输入同时带相反标签，硬间隔约束将不可行，应重新考虑软间隔。}

\problem{Q33}{拉格朗日、对偶与KKT：把求导串起来}{背面第1栏／S02，必要时S01}
\qtext{对线性硬间隔SVM，写出Lagrangian，对$w,b$求导并推出对偶目标与约束。然后用$(-1,-1),(1,+1)$验证对偶最优解，说明怎样由乘子恢复$w,b$，并检查原始与对偶目标值。}
\route{“Lagrangian、求导、dual、KKT”指向S02。先写清原始变量与乘子，再消元；一般矩阵求导只负责代数步骤，SVM的标签约束决定了具体符号。}
\given{一般情形$x_i,w\in\R^d$，$y_i=\pm1$，$b\in\R$。每个不等式$y_i(w\T x_i+b)-1\ge0$对应$\alpha_i\ge0$。}
\solve
构造并明确鞍点方向：
\[
L(w,b,\alpha)=\frac12w\T w-\sum_i\alpha_i\{y_i(w\T x_i+b)-1\},
\qquad \min_{w,b}\max_{\alpha\ge0}L.
\]
对原始变量求偏导：
\[
\frac{\partial L}{\partial w}=w-\sum_i\alpha_i y_i x_i=0,
\qquad
\frac{\partial L}{\partial b}=-\sum_i\alpha_i y_i=0.
\]
所以$w=\sum_i\alpha_i y_i x_i$、$\sum_i\alpha_i y_i=0$。代回得到
\[
\max_{\alpha}\ D(\alpha)=\sum_i\alpha_i-
\frac12\sum_{i,j}\alpha_i\alpha_jy_iy_jx_i\T x_j,
\quad \alpha_i\ge0,\quad\sum_i\alpha_i y_i=0.
\]
KKT还需原始可行及$\alpha_i\{y_i g_i-1\}=0$。
对本题两个样本，和约束给$\alpha_1=\alpha_2=a$，且$w=2a$，故
\[
D(a)=2a-2a^2,\quad a\ge0;
\qquad D'(a)=2-4a=0\Rightarrow a=\frac12.
\]
$D''=-4<0$，这是凹二次函数的全局最大值。由$w=2a$得$w=1$；因$\alpha_2>0$，$g(1)=1$，所以$b=0$。

\ans{最优乘子为$(1/2,1/2)$；恢复$(w,b)=(1,0)$，原始最小值与对偶最大值均为$1/2$。对偶是最大化，L对$w,b$是最小化方向。}
\qcheck{非负乘子、$\alpha\T y=0$、两个训练约束、互补条件全部成立；$\tfrac12w^2=D(1/2)=1/2$。}
\pitfall{不要把L同时对$w,b,\alpha$最小化；不要把普通的$\sum_i\alpha_i=0$误写成分类约束，它必须包含$y_i$。}
\transfer{换成核时，只把$x_i\T x_j$替换为$K(x_i,x_j)$；若改软间隔，还要加入$\alpha_i\le C$，其他标签符号不变。}

\problem{Q34}{支持向量预测，以及“在边界上”的例外}{背面第1栏／S02}
\qtext{一维硬间隔训练集按顺序为
$x=(-2,-1,1,1,2)$，$y=(-1,-1,+1,+1,+1)$，给定一组最优乘子
$\alpha=(0,1/2,1/2,0,0)$、$b=0$。求$w$与$g(-1/4)$；列出各训练点的函数间隔。判断：“只要样本在间隔边界上，它的$\alpha$必定大于0。”再检查有人写的$g(x)=\sum_i\alpha_i x_i x+b$。}
\route{“support vector、给定乘子、预测、边界”对应S02。先恢复带标签的$w$，再用KKT分清单向条件；重复样本可暴露边界点与非零乘子的区别。}
\given{样本3与4输入及标签相同，是两个训练条目。支持向量展开中采用非零乘子定义；不把所有几何边界点都强行计为非零系数。}
\solve
正确恢复式包含标签：
\[
w=\sum_i\alpha_i y_i x_i
=\frac12(-1)(-1)+\frac12(+1)(1)=1.
\]
因此$g(x)=x$，$g(-1/4)=-1/4<0$，预测$-1$。
\[
(y_1g_1,\ldots,y_5g_5)=(2,1,1,1,2).
\]
点2、3、4都在规范间隔边界上，但点4的$\alpha_4=0$。这是允许的退化情形：其重复点3已承担相应系数。
互补条件为$\alpha_i(m_i-1)=0$，它只能推出
\[
\alpha_i>0\Rightarrow m_i=1,
\qquad m_i>1\Rightarrow\alpha_i=0.
\]
$m_i=1$时，乘积对任何非负$\alpha_i$都等于0，不能单凭这一式反推$\alpha_i>0$。

错误预测式漏$y_i$，在本题得到$\sum_i\alpha_i x_i=-1/2+1/2=0$，把所有输入都预测为边界值0。应改为
\[
g(x)=\sum_i\alpha_i y_i\,x_i x+b.
\]
\ans{$w=1$、$g(-1/4)=-1/4$，预测$-1$。在间隔边界上不保证该条目的乘子非零；分类预测展开必须保留$y_i$。}
\qcheck{$\sum\alpha_i y_i=0$，所有$m_i\ge1$，$\alpha_i(m_i-1)=0$；原始与对偶目标均为$1/2$。}
\pitfall{支持向量条件与几何边界条件不可随意写成双向等价。若$g=0$，类别必须按题目指定的零值约定处理，不能擅自给出唯一类别。}
\transfer{核分类中将$x_i x$换成$K(x_i,x)$，标签仍保留。SVR与本册LS-SVM回归使用的展开系数有不同定义，不能机械多乘一个$y_i$。}

\problem{Q35}{软间隔的松弛：正确也可能有惩罚}{背面第2栏／S03}
\qtext{软间隔SVM使用$y_i g_i\ge1-\xi_i$、$\xi_i\ge0$，$C>0$。五个样本的函数间隔依次为$m=(1.4,1,0.6,0,-0.4)$。求固定$w,b$时最小可行松弛，并分别说明它们是间隔外、间隔边界、正确但在间隔内、决策面或错分。若$\norm w=2,C=2$，求此时目标值。}
\route{圈出“$\xi$、margin violation、$C$”。S03用$y_i g_i$判断分类与间隔；ADALINE的线性残差和SVR的$\varepsilon$管道均不采用这一标签乘积。}
\given{$m_i=y_i g_i$；在固定模型下，目标含正惩罚$C\sum\xi_i$，所以应取每个样本最小可行$\xi_i=\max(0,1-m_i)$。}
\solve
约束逐项改写为$\xi_i\ge1-m_i$，再与$\xi_i\ge0$合并：
\[
\xi=(0,0,0.4,1,1.4),\qquad\sum_i\xi_i=2.8.
\]
\begin{center}\begin{tabular}{@{}ccl@{}}\toprule
$m_i$&$\xi_i$&状态\\\midrule
1.4&0&正确、在规范间隔外\\
1&0&正确、在规范间隔边界\\
0.6&0.4&正确、但在间隔内，存在间隔违背\\
0&1&在决策面，最终归类取决于零值约定\\
$-0.4$&1.4&错分，因$y_i g_i<0$\\\bottomrule
\end{tabular}\end{center}
代入目标时不要漏范数平方：
\[
J=\frac12\norm w^2+C\sum_i\xi_i
=\frac12(2)^2+2(2.8)=7.6.
\]
只有在取最小松弛的约定下，$0<\xi<1$表示正确但在间隔内，$\xi>1$表示错分；可行但人为放大的松弛不能直接用于反推样本位置。

\ans{最小松弛为$(0,0,0.4,1,1.4)$，固定模型的目标值为$7.6$。间隔违背不等于错分；$\xi=1$对应决策面。}
\qcheck{逐项算$m_i+\xi_i=(1.4,1,1,1,1)$，均不小于1；松弛均非负。}
\pitfall{不要用$|y_i-g_i|$替代$\max(0,1-y_i g_i)$；软间隔惩罚的是间隔违背，包含一部分仍被正确分类的点。}
\transfer{若只改变$C$且固定$w,b$，松弛值不变，目标的惩罚权重改变；若重新训练，$w,b$、各$m_i$和$\xi_i$都应重算。}

\problem{Q36}{C、盒约束与KKT状态：哪些能反推？}{背面第2栏／S03，背面第1栏／S02}
\qtext{软间隔SVM取$C=2$。某候选解的三类条目为：A的$(\alpha,m,\xi)=(0,1.5,0)$；B为$(1,1,0)$；C为$(2,0.4,0.6)$。检查各自KKT互补条件。判断“$\alpha=C$必错分”和“提高$C$必使测试误差降低”。然后对$(-1,-1),(1,+1)$，改取$C=1/4$，求最优乘子、$w$及一个可用$b$。}
\route{“$0\le\alpha\le C$、free support vector、参数$C$”对应S03。盒约束来自线性松弛惩罚；本册LS-SVM的等式乘子没有这个盒。}
\given{软间隔KKT除可行性外含
$\alpha_i(m_i-1+\xi_i)=0$、$(C-\alpha_i)\xi_i=0$。
全局还需$w=\sum\alpha_i y_i x_i$和$\sum\alpha_i y_i=0$。}
\solve
A的两乘积均为0，且$m>1,\xi=0$；B有$0<\alpha<C$，于是必须$\xi=0,m=1$；C有$m-1+\xi=0$且$C-\alpha=0$，也满足局部互补条件。仅凭这三项仍不能确认整组数据的全局驻点关系。

一般结论为
\[
\alpha=0\Rightarrow m\ge1,\quad
0<\alpha<C\Rightarrow m=1,\quad
\alpha=C\Rightarrow m\le1.
\]
最后一种可处在边界、正确但在间隔内、决策面或错分；C条目$m=0.4>0$即为反例。
提高$C$使训练中的间隔违背代价更大，通常推动更紧地适配训练数据；泛化仍取决于数据、核及参数选择，不能保证测试误差降低。

两样本的小题依然有$\alpha_1=\alpha_2=a$、$D(a)=2a-2a^2$，但现在
\[
0\le a\le\frac14.
\]
$D'(a)=2-4a>0$在该区间成立，所以$a=1/4$；$w=2a=1/2$。取$b=0$，两点都有$m=1/2,\xi=1/2$。
\[
J=\frac12(1/2)^2+\frac14(1/2+1/2)=\frac38=D(1/4).
\]
\ans{A、B、C的局部互补条件均可成立；$\alpha=C$不等于错分。$C=1/4$的两样本解可取$\alpha=(1/4,1/4),w=1/2,b=0$、目标$3/8$。}
\qcheck{两乘子均在盒内，$\sum\alpha y=0$；两个原始约束取等号；原始与对偶目标相等。此题$b$不唯一，$[-1/2,1/2]$都可行且最优。}
\pitfall{不要在$\alpha=C$时无条件用$yg=1$求$b$；该等式用于$0<\alpha<C$的自由支持向量。}
\transfer{若$C\ge1/2$，本数据的无盒约束最优$a=1/2$已允许，恢复硬间隔解。换数据后阈值应重新求，不能普遍套$C\ge1/2$。}

\problem{Q37}{多项式核：亲手展开六维映射}{背面第2栏／K01}
\qtext{二维输入$x=(x_1,x_2)\T$采用$K(x,z)=(1+x\T z)^2$。构造一个显式特征映射$\phi(x)$使$K(x,z)=\phi(x)\T\phi(z)$。对$a=(1,0)\T,b=(0,1)\T$，分别用核公式与特征内积算$K(a,b)$、$K(a,a)$。}
\route{“kernel trick、explicit feature map、多项式展开”查K01。不是先训练网络隐藏权重；核在此定义了内积，不需要为每个特征分量另求训练参数。}
\given{原输入维度2，展开后可取6维。核内积中的交叉项系数必须由特征的缩放因子共同产生。为避免与偏置混淆，本题$b$仅为第二个输入向量的名称。}
\solve
令$u=x_1z_1,v=x_2z_2$，则
\[
(1+u+v)^2=1+u^2+v^2+2uv+2u+2v.
\]
所以可选择
\[
\phi(x)=\begin{bmatrix}
1&x_1^2&\sqrt2x_1x_2&x_2^2&\sqrt2x_1&\sqrt2x_2
\end{bmatrix}\T\in\R^6.
\]
对应内积逐项为
\[
1+x_1^2z_1^2+2x_1x_2z_1z_2+x_2^2z_2^2+2x_1z_1+2x_2z_2,
\]
恰等于核展开。特征顺序可改变，只要两个向量同步使用相同顺序。

对给定输入，$a\T b=0,a\T a=1$，直接计算
\[
K(a,b)=1,\qquad K(a,a)=4.
\]
另一条路线为
\[
\phi(a)=(1,1,0,0,\sqrt2,0)\T,\quad
\phi(b)=(1,0,0,1,0,\sqrt2)\T.
\]
故$\phi(a)\T\phi(b)=1$，而$\norm{\phi(a)}^2=1+1+2=4$。

\ans{上述六维$\phi$给出所需内积；$K(a,b)=1,K(a,a)=4$。核技巧允许直接算核值而无需实际存储六维向量。}
\qcheck{检查平方交叉项是否为2倍；对角核必须等于特征范数平方，因此本题$K(a,a)$不可能为负。}
\pitfall{漏掉$\sqrt2$会把对应的2倍项变成1倍；$(x\T z)^2$是齐次核，与本题$(1+x\T z)^2$的常数及线性特征不同。}
\transfer{改核次数或常数后，重新展开并核验系数；不直接沿用六维映射。预测依然是在特征空间做仿射函数，在原输入中一般非线性。}

\problem{Q38}{Gram矩阵与无效核：对称还不够}{背面第2栏／K01}
\qtext{对$a=(1,0)\T,b=(0,1)\T$使用$K(x,z)=(1+x\T z)^2$，构造$2\times2$的Gram矩阵并判断半正定。另一个候选“核”在某两个样本上产生矩阵
$G=\bigl[\begin{smallmatrix}1&2\\2&1\end{smallmatrix}\bigr]$。有人认为“对称且所有元素为正，所以是有效核”，请给出反例计算。最后说明一个数据集的Gram半正定能否证明核在全域有效。}
\route{“Gram、PSD、有效kernel”对应K01。核有效性检查是内积几何的条件；LS-SVM中的$\gamma^{-1}I$是正则项，不能用它掩盖原核定义无效。}
\given{Gram第$(i,j)$项为$K(x_i,x_j)$，矩阵必须对称且对任意$c$满足$c\T Kc\ge0$。有效核要求对\textbf{所有有限样本集合}都满足该条件。}
\solve
依次计算对角与交叉项：
\[
K(a,a)=4,\quad K(a,b)=K(b,a)=1,\quad K(b,b)=4,
\qquad K=\begin{bmatrix}4&1\\1&4\end{bmatrix}.
\]
其特征值为$4+1=5$、$4-1=3$，均为正；也可对任意$c=(c_1,c_2)\T$写
\[
c\T Kc=4c_1^2+2c_1c_2+4c_2^2
=3(c_1^2+c_2^2)+(c_1+c_2)^2\ge0.
\]
对候选$G$，取$c=(1,-1)\T$即可否定：
\[
Gc=(-1,1)\T,\qquad c\T Gc=-2<0.
\]
尽管$G$对称且逐元素为正，它不是半正定矩阵，因而不能是实特征内积在这两点上的Gram矩阵。

一个集合检查通过，仅说明该集合未发现违反条件，不能证明全域有效。本题的多项式核可借Q37的显式映射作全域证明：
$c\T Kc=\norm{\sum_i c_i\phi(x_i)}^2\ge0$对任何有限样本都成立。

\ans{多项式核Gram为$\bigl[\begin{smallmatrix}4&1\\1&4\end{smallmatrix}\bigr]$，正定；候选$G$因$c=(1,-1)\T$给出二次型$-2$而无效。逐元素正不等于矩阵半正定。}
\qcheck{Gram的对角项应为$\norm{\phi(x_i)}^2$，矩阵与其转置相等；PSD二次型测试不能仅代$c=(1,1)\T$。}
\pitfall{不能把“对称”或“一组样本上PSD”单独当成全域核有效的充分条件；sigmoid表达式也不是任意参数下都可直接认定为有效核。}
\transfer{换样本时逐项重建Gram；对小对称矩阵可用特征值或二次型。维度变大时原理不变，不能只检查元素符号。}

\problem{Q39}{Gaussian核预测：标签与宽度分母都要查}{背面第2栏／K01，背面第1栏／S02}
\qtext{二类训练输入$x_1=-1,x_2=1$，标签$y=(-1,+1)$，使用
$K(x,z)=\exp[-(x-z)^2/(2\sigma^2)]$、$\sigma=1$。给定软间隔解$\alpha_1=\alpha_2=C=1,b=0$。构造Gram矩阵，求$g(1)$及类别，并检查这组解的松弛和KKT。若误把分母写成$\sigma^2$，会得到什么不同的核值与分数？}
\route{“Gaussian/RBF kernel、给定$\alpha$、分类预测”查K01与S02。先核对分类标签，再核对宽度约定；此处是核SVM，不能拿RBF网络的任意输出权重替代$\alpha_i y_i$。}
\given{距离平方$(x_1-x_2)^2=4$。令$\rho=e^{-2}\approx0.135335$，则Gram非对角项为$\rho$。}
\solve
\[
K=\begin{bmatrix}1&\rho\\\rho&1\end{bmatrix},\qquad
g(z)=\sum_i\alpha_i y_iK(x_i,z)+b=-K(-1,z)+K(1,z).
\]
在$z=1$处，
\[
g(1)=1-e^{-2}\approx0.864665>0,
\]
故类别为$+1$。两训练点的函数间隔都为$1-\rho$，最小松弛均为$\rho$。它们仍被正确分类，但在间隔内。
两乘子在盒的上界，$\sum_i\alpha_i y_i=0$；$m_i-1+\xi_i=0$，且$(C-\alpha_i)\xi_i=0$。核空间范数平方为
\[
\norm w^2=\sum_{i,j}\alpha_i\alpha_jy_iy_jK_{ij}=2(1-\rho).
\]
原始目标为$(1-\rho)+2\rho=1+\rho$，对偶值为$2-(1-\rho)=1+\rho$，相等。

若错误改分母为$\sigma^2$而仍声称$\sigma=1$不变，交叉核变为$e^{-4}\approx0.018316$，分数变为$1-e^{-4}\approx0.981684$。这实际上改了核宽度，不能视为同一数值计算。

\ans{按题目约定，Gram为$\bigl[\begin{smallmatrix}1&e^{-2}\\e^{-2}&1\end{smallmatrix}\bigr]$，$g(1)=1-e^{-2}>0$判$+1$。宽度分母改变后须重新解释参数。}
\qcheck{对角核都为1；$g(-1)=-g(1)$；原始与对偶均为$1+e^{-2}$。}
\pitfall{分类核展开漏$y_i$会把两项都加正；RBF核相似度是非负数，但判别分数可以为负。}
\transfer{若输入缩放或改变$\sigma$，重算距离与核；较小$\sigma$使相似度更局部。选择$C,\sigma$应使用验证资料，不能保证越小$\sigma$泛化越好。}

\problem{Q40}{两组比较：相同Gaussian不等于相同模型}{背面第2栏／D05；RBF细节见正面第5栏／R01}
\qtext{用完整句子比较：(a) Hard-margin V.S. Soft-margin SVM；(b) RBF network V.S. RBF-kernel SVM。每组写共同点、明确维度的差异、适用任务及限制。说明为什么“都使用Gaussian，所以训练方法与中心选择完全相同”不正确。}
\route{“compare、共同点、区别”先查D05；见Gaussian时继续辨认是网络隐层还是核分类器。仅凭函数名称不能完成模型识别。}
\given{本题RBF网络为$f(x)=\sum_{j=1}^{M}v_j h_j(x)+b$，$h_j$由中心$c_j$与宽度控制；核SVM分类为$g(x)=\sum_i\alpha_i y_iK(x_i,x)+b$。}
\solve
\textbf{(a) 可写成答卷的比较。}两者都用范数惩罚控制分类面的复杂度，并依据带标签间隔训练；都可以在特征空间构造分类面。硬间隔要求所有训练点满足$y_i g_i\ge1$，所以必须在选定特征空间可分；它不允许间隔违背。软间隔加入$\xi_i\ge0$，满足$y_i g_i\ge1-\xi_i$，并最小化$\tfrac12\norm w^2+C\sum\xi_i$，适合有噪声、重叠或不能完全分开的数据。软间隔允许一部分训练误差，$C$控制惩罚强度。硬间隔可能被离群点拉动，软间隔也需选择$C$和核参数，均不保证任何测试集上最优。

\textbf{(b) 可写成答卷的比较。}两者都可用Gaussian相似度表示输入的局部关系，最终输出都可写成若干基函数的加权和。RBF网络显式设置有限个隐藏单元，中心与宽度可以预选或通过数据学习，再训练输出权重；可用于回归或分类。RBF-kernel SVM通过核内积和间隔目标解受约束二次优化；核中心在展开中对应训练样本，只有非零乘子的支持向量进入分类预测，系数为$\alpha_i y_i$且受训练约束。它通常用于分类；同族SVR处理回归。RBF网络的任意输出权重不受同一套SVM盒约束；RBF-kernel SVM也不是先独立训练一组可自由移动的隐藏中心。

\ans{共同Gaussian给出相似的基函数形式；隐藏中心/系数的来源、损失与约束、训练方法不同，所以不能认定是同一个模型。关键词：maximum margin、slack、box constraint、center、basis function、support vector。}
\qcheck{每组都应明确“按哪个维度比较”，而且同时写共同点、任务与限制；不只写“都是机器学习”。}
\pitfall{BP是计算梯度的训练方法，不是与RBF或SVM并列的同类模型；同一个公式形式也可能来自不同训练原则。}
\transfer{题目换成预测数值时，先检查使用平方误差、$\varepsilon$不敏感损失还是其他目标；RBF这一名称本身不能决定损失和训练算法。}

\problem{Q41}{SVR：由连续输出与残差建立两侧管道}{背面第3栏／V01}
\qtext{希望用$f(x)=w\T\phi(x)+b$预测连续值，允许绝对残差在$\varepsilon=1/2$以内不计损失，超出部分线性惩罚，同时控制$\norm w$。写出目标、两类松弛与约束。固定模型后，四点残差$r_i=y_i-f_i$依次为$0,0.3,1.2,-0.9$，计算$\xi_i,\xi_i^*$及总$\varepsilon$不敏感损失。}
\route{圈出“continuous output、$\varepsilon$、tube、超出线性惩罚”。V01对应SVR；分类SVM的$y_i g_i$用于$\pm1$标签，LS-SVM则使用平方误差与等式约束。}
\given{$r=y-f$是有符号残差，$\xi$负责预测偏低的一侧$r>\varepsilon$，$\xi^*$负责预测偏高的一侧$r<-\varepsilon$。$C>0,\varepsilon\ge0$。}
\solve
写模型后构建优化问题：
\[
\min_{w,b,\xi,\xi^*}\frac12\norm w^2+C\sum_i(\xi_i+\xi_i^*),
\qquad
\begin{cases}
y_i-f_i\le\varepsilon+\xi_i,\\
f_i-y_i\le\varepsilon+\xi_i^*,\\
\xi_i,\xi_i^*\ge0.
\end{cases}
\]
固定模型的最小可行松弛为
\[
\xi_i=\max(0,r_i-\varepsilon),\quad
\xi_i^*=\max(0,-r_i-\varepsilon),\quad
\ell_\varepsilon(r_i)=\max(0,|r_i|-\varepsilon).
\]
\begin{center}\begin{tabular}{@{}cccc@{}}\toprule
$r_i$&$\xi_i$&$\xi_i^*$&$\ell_\varepsilon(r_i)$\\\midrule
0&0&0&0\\
0.3&0&0&0\\
1.2&0.7&0&0.7\\
$-0.9$&0&0.4&0.4\\\bottomrule
\end{tabular}\end{center}
总损失为$1.1$，松弛惩罚为$1.1C$，还须加$\tfrac12\norm w^2$才是总目标。管道的全高度是$2\varepsilon=1$；这里按输出轴的误差定义管道，不是把二维图上任意斜率的垂直距离直接当几何正交距离。

\ans{SVR在$|r|\le1/2$时不惩罚，外部线性惩罚；本题$\xi=(0,0,0.7,0)$、$\xi^*=(0,0,0,0.4)$，总损失$1.1$。}
\qcheck{逐项检查$r_i\le1/2+\xi_i$及$-r_i\le1/2+\xi_i^*$；两侧绝不能漏一条。}
\pitfall{本册统一$f=w\T\phi+b$，残差是$y-f$；不要混用减$b$的另一个偏置约定。$\xi^*$不是$\xi$的复共轭，而是另一侧松弛。}
\transfer{换$\varepsilon$要重新判断哪些点位于管道内；若只固定模型改$C$，松弛不变，但重新训练后模型和残差都可能改变。}

\problem{Q42}{SVR的对偶：为什么是差系数？}{背面第3栏／V01、V02}
\qtext{沿用$f_i=w\T\phi_i+b$、$r_i=y_i-f_i$与两侧SVR约束，给第一侧乘子$\alpha_i\ge0$、第二侧$\alpha_i^*\ge0$，非负松弛乘子为$\eta_i,\eta_i^*\ge0$。构建Lagrangian，求导并写出核对偶及预测式；解释为什么预测用$\alpha_i-\alpha_i^*$，不能用分类SVM的$\alpha_i y_i$。}
\route{“两侧约束、dual、Lagrangian、差系数”查V01；恢复偏置与边界条件查V02。不要只因为模型名称含SVM，就移植分类标签乘子。}
\given{$\phi_i=\phi(x_i)$，$K_{ij}=\phi_i\T\phi_j$；两组$\alpha$对应两个不等式，$y_i\in\R$是连续目标。令$d_i=\alpha_i-\alpha_i^*$。}
\solve
\begin{align*}
L={}&\frac12\norm w^2+C\sum_i(\xi_i+\xi_i^*)-\sum_i(\eta_i\xi_i+\eta_i^*\xi_i^*)\\
&-\sum_i\alpha_i(\varepsilon+\xi_i-y_i+w\T\phi_i+b)\\
&-\sum_i\alpha_i^*(\varepsilon+\xi_i^*+y_i-w\T\phi_i-b).
\end{align*}
原始变量取最小、非负乘子取最大。驻点条件为
\[
w-\sum_i(\alpha_i-\alpha_i^*)\phi_i=0,\quad
\sum_i(\alpha_i^*-\alpha_i)=0,
\quad C-\alpha_i^{(*)}-\eta_i^{(*)}=0.
\]
所以$w=\sum_i d_i\phi_i$、$\sum_i d_i=0$，并由乘子非负得到$0\le\alpha_i,\alpha_i^*\le C$。消元后
\[
\max_{\alpha,\alpha^*}\left\{
-\frac12d\T Kd-\varepsilon\sum_i(\alpha_i+\alpha_i^*)+y\T d\right\},
\quad\mathbf1\T d=0,\quad0\le\alpha_i,\alpha_i^*\le C.
\]
预测为$f(x)=\sum_i d_iK(x_i,x)+b$。两侧误差作用相反，因此形成差系数；连续$y_i$只在对偶的$y\T d$里，不额外乘进预测系数。

两组互补条件还必须保留：
\[
\alpha_i(\varepsilon+\xi_i-r_i)=0,\quad
\alpha_i^*(\varepsilon+\xi_i^*+r_i)=0,
\quad(C-\alpha_i)\xi_i=(C-\alpha_i^*)\xi_i^*=0.
\]
\ans{SVR预测系数是$d=\alpha-\alpha^*$，满足和为0；对偶的$\varepsilon$项用两组乘子的\textbf{和}，二次项与目标数据项用两组乘子的\textbf{差}。}
\qcheck{从$\partial L/\partial b=0$检查差系数之和；从$\partial L/\partial w$检查预测符号，不能仅凭记忆调换两组乘子。}
\pitfall{把$-\varepsilon\sum(\alpha+\alpha^*)$误写为差会改变优化问题；同样不能漏掉任一组盒约束。}
\transfer{换核只替换$K$；若更换残差的符号定义，必须同步交换两类松弛与乘子，不能只改预测一处。}

\problem{Q43}{自由支持向量恢复偏置：完整SVR小例题}{背面第3栏／V01、V02}
\qtext{两点为$(x_1,y_1)=(0,0),(x_2,y_2)=(2,2)$，线性核$K(x,z)=xz$，$\varepsilon=1/2,C=1$。已知$\alpha=(0,1/4)$、$\alpha^*=(1/4,0)$。恢复$w,b$并预测$x=1$；求训练残差和两组松弛，逐项核验KKT及原始、对偶目标。}
\route{“$0<\alpha<C$、恢复$b$、SVR”查V02。SVR的边界是$r=\pm\varepsilon$；分类SVM的边界$yg=1$不适用于连续目标。}
\given{$d=(-1/4,1/4)\T$，$K=\bigl[\begin{smallmatrix}0&0\\0&4\end{smallmatrix}\bigr]$。写$q_i=\sum_jd_jK(x_j,x_i)$，于是$f_i=q_i+b$。}
\solve
先恢复线性系数与不含偏置的训练分数：
\[
w=\sum_i d_ix_i=\frac12,\qquad q=(0,1)\T.
\]
点2有$0<\alpha_2<C$，所以$\xi_2=0,r_2=\varepsilon$，可用
\[
b=y_2-q_2-\varepsilon=2-1-\frac12=\frac12.
\]
点1有$0<\alpha_1^*<C$，所以$\xi_1^*=0,r_1=-\varepsilon$，独立给出
\[
b=y_1-q_1+\varepsilon=0-0+\frac12=\frac12.
\]
故$f(x)=x/2+1/2$，$f(1)=1$；训练预测$(1/2,3/2)$，残差$r=(-1/2,1/2)$，两类松弛全部为0。

检查乘子均在$[0,1]$、$\sum d=0$、$w=\sum d_ix_i$。非零$\alpha_2$所乘的管道约束在上侧取等号，非零$\alpha_1^*$所乘的约束在下侧取等号；其他乘子为0，所有互补乘积均为0。
\[
J=\frac12(1/2)^2=\frac18,
\]
\[
D=-\frac12\left(\frac14\right)-\frac12\left(\frac12\right)
+\left[0\left(-\frac14\right)+2\left(\frac14\right)\right]
=\frac18.
\]
这里$d\T Kd=1/4$，$\sum(\alpha+\alpha^*)=1/2$。

\ans{$w=b=1/2$，$f(1)=1$；训练残差为$(-1/2,1/2)$，零松弛，原始和对偶目标均为$1/8$。}
\qcheck{两类自由SV各给相同$b$；把$b$代回检查残差的正负比单独记一个偏置公式可靠。}
\pitfall{$\alpha$侧给$b=y-q-\varepsilon$，$\alpha^*$侧给$b=y-q+\varepsilon$，不能交换。管道边界可有非零系数，不能写成“只有管道外才是SV”。}
\transfer{如果改变数据、核、$C$或$\varepsilon$，原乘子不一定仍最优。必须先重新确认$0<\alpha^{(*)}<C$，才沿用自由SV偏置公式。}

\problem{Q44}{无自由SV的偏置区间与零管宽例外}{背面第3栏／V02}
\qtext{(a) 两点$(0,0),(1,2)$，线性核，$C=1/2,\varepsilon=1/2$；给定$\alpha=(0,1/2)$、$\alpha^*=(1/2,0)$。求$w$和所有可行最优$b$，选$b=3/4$检查残差、松弛和目标。(b) 当$\varepsilon=0$时，“同一点$\alpha_i,\alpha_i^*$绝不能同时正”是否仍成立？用单点$(x,y)=(0,2)$、$C=1$构造或否定一个反例。}
\route{“no free support vector、乘子在盒端点、$\varepsilon=0$”查V02。此题没有$0<\alpha^{(*)}<C$，所以不能无条件把残差设为管道边界。}
\given{(a)$d=(-1/2,1/2)$，$w=\sum d_ix_i=1/2$、$q=(0,1/2)$。四条KKT偏置规则分别来自两侧约束与松弛互补。}
\solve
在没有自由系数时收集以下界：
\[
\begin{array}{ll}
\alpha_i=0:\ b\ge y_i-q_i-\varepsilon,&
\alpha_i=C:\ b\le y_i-q_i-\varepsilon,\\
\alpha_i^*=0:\ b\le y_i-q_i+\varepsilon,&
\alpha_i^*=C:\ b\ge y_i-q_i+\varepsilon.
\end{array}
\]
点1给$b\ge-1/2$和$b\ge1/2$；点2给$b\le1$和$b\le2$，交集为
\[
\boxed{\frac12\le b\le1}.
\]
取$b=3/4$，则$f=(3/4,5/4)$，$r=(-3/4,3/4)$。
$\xi=(0,1/4)$、$\xi^*=(1/4,0)$；正松弛对应的乘子均在上界$C$，所乘约束取等号。
\[
J=\frac12(1/2)^2+\frac12(1/4+1/4)=\frac38,
\]
\[
D=-\frac12(1/4)-\frac12(1)+2(1/2)=\frac38.
\]
所以区间内的$b$都能完成同一最优$d$对应的原始解；其中端点也须按KKT检查，不能只用某个“平均值公式”。

\textbf{(b)}若$\varepsilon>0$且两乘子都正，两个互补等式相加会要求
$2\varepsilon+\xi+\xi^*=0$，与非负松弛冲突。
但$\varepsilon=0$时这一排除理由失效。对给定单点取
\[
w=0,\ b=2,\ \xi=\xi^*=0,\ \alpha=\alpha^*=1/4.
\]
差系数为0，和约束、驻点与所有互补条件均成立，原始及对偶值均为0；两乘子确实可同时为正。
\ans{(a)$w=1/2,b\in[1/2,1]$，所选$b=3/4$给目标$3/8$。(b)“两乘子不能同时正”需$\varepsilon>0$；$\varepsilon=0$应重新检查，预测始终使用差系数。}
\qcheck{不要只验$\sum d=0$；本题通过原始约束、互补和目标相等共同验证。}
\pitfall{把所有非零乘子都当自由SV，会在(a)误得唯一$b$；在(b)忽略$\varepsilon=0$会误删合法的对偶解。}
\transfer{换数据后仍逐项收集$b$的下界与上界；若最大下界大于最小上界，说明候选系数或数值容差存在问题，应重新检查，不能随意平均。}

\problem{Q45}{SVM、SVR、LS-SVM：用任务、损失和约束辨认}{背面第5栏／D06；S01、V01、L01辅助}
\qtext{辨认并比较三种训练描述：(A)输出二类标签$\pm1$，约束$y_i g_i\ge1-\xi_i$，惩罚$C\sum\xi_i$；(B)预测实数，管道内不罚，外部线性惩罚；(C)预测实数，$y_i=f_i+e_i$，目标含$\gamma\sum e_i^2/2$。写出模型名称、共同点、关键差异与限制。判断：“SVM只能分类，LS-SVM只能回归”“求解线性方程组意味着目标不是二次优化”。}
\route{先圈出标签类型和输出，再检查损失及约束。A的标签间隔指向SVM；B的$\varepsilon$管道指向SVR；C的平方误差加等式指向本课件的LS-SVM回归。仅出现$\gamma$或kernel一词不足以确定模型。}
\given{三者均可用$f=w\T\phi+b$并通过核求内积；分类时再对判别分数按符号给标签。比较“标准稀疏性”指常见最优展开性质，不是每个数据集的绝对保证。}
\solve
\textbf{完整句子答案。}A是软间隔支持向量分类机，采用hinge形式的间隔违背惩罚；B是$\varepsilon$不敏感支持向量回归SVR；C是最小二乘支持向量机（Least-Squares Support Vector Machine，LS-SVM）的回归形式。三者共同用范数惩罚控制函数复杂度，核技巧可表达非线性关系；它们都要选择正则及核参数并评估泛化。

三者在明确维度上的差异如下：
\begin{center}\small\begin{tabular}{@{}p{24mm}p{40mm}p{40mm}p{44mm}@{}}\toprule
维度&SVM分类&SVR&本课LS-SVM回归\\\midrule
输出与损失&类别；$\max(0,1-yg)$&实数；$\max(0,|y-f|-\varepsilon)$&实数；平方残差$e^2$\\
约束&标签间隔不等式&上下管道不等式&$y=f+e$等式\\
求解&受约束二次优化&受约束二次优化&二次优化条件化为线性系统\\
展开系数&$\alpha_i y_i$&$\alpha_i-\alpha_i^*$&$\alpha_i=\gamma e_i$\\
乘子范围&$0\le\alpha_i\le C$&两组均在$[0,C]$&等式乘子可正可负\\
展开性质&通常稀疏&通常稀疏&通常稠密\\\bottomrule
\end{tabular}\end{center}
两句判断都错误。SVM方法族包括分类与SVR回归，LS-SVM也可用于分类；本课件只对回归给出完整训练推导，分类应用图不等于只能做回归。LS-SVM原目标仍是带等式约束的二次优化，求导消元后能直接解线性方程组。

\ans{A＝SVM分类，B＝SVR，C＝LS-SVM回归。共同点应写“范数正则、核内积、模型选择”；差异应落到输出、损失、约束、系数与训练求解上。}
\qcheck{检查是否误写“等式要求穿过全部样本”：C中$e$允许非零；是否误把LS-SVM乘子也限制为非负或盒约束。}
\pitfall{目标含平方项不代表只能做线性预测；线性系统描述未知参数的求解，核仍可使输入到输出非线性。}
\transfer{遇到新任务先问输出是类别还是数值，再据损失、约束和参数条件选择方法。模型选择理由应联系噪声、误差容忍与可用验证资料，而非只凭模型名称。}

\problem{Q46}{从电负荷预测文字构建LS-SVM回归}{背面第4栏 L01、L02；正面第1栏 M02}
\qtext{有$n\ge1$个训练样本$(x_i,y_i)$，其中$x_i\in\R^d$为预测前已经可用的特征，$y_i\in\R$为实际电负荷。指定有效核$K(x,z)=\phi(x)^T\phi(z)$、$\gamma>0$，要求预测包含常数偏置，并同时控制权重范数和训练平方误差。请自行写出预测函数、目标、约束、Lagrangian、四组导数及最终线性系统。采用未知顺序$[\alpha;b]$。解释等式是否强制插值，以及模型最终输出什么。}
\route{圈“连续电负荷、平方误差、范数、偏置、有效核、$\gamma$”。这些条件共同指向LS-SVM回归；只有“least squares”一词并不足以区分LMS或固定RBF输出拟合。这里指定核和带范数惩罚的等式建模，使用L01；不是分类SVM的$y_i=\pm1$间隔。}
\given{$\phi_i=\phi(x_i)\in\R^h$；未知$w\in\R^h,b\in\R,e,\alpha\in\R^n$；$K\in\R^{n\times n}$。$e_i$定义为实际减预测，$\gamma$为正的误差权重。}
\solve
先写$f(x)=w^T\phi(x)+b$，由“实际=预测+残差”得到
\[
\min_{w,b,e}\frac12w^Tw+\frac\gamma2\sum_{i=1}^ne_i^2,
\qquad y_i=w^T\phi_i+b+e_i.
\]
范数项限制复杂度，平方项衡量拟合；$1/2$让求导时的2消去。等式允许$e_i\ne0$，所以不强制通过全部样本。
\[
\mathcal L=\frac12w^Tw+\frac\gamma2\sum_i e_i^2-\sum_i\alpha_i(w^T\phi_i+b+e_i-y_i).
\]
逐组令偏导为零：
\begin{align*}
\partial_w\mathcal L=0&\Rightarrow w=\sum_i\alpha_i\phi_i, &
\partial_b\mathcal L=0&\Rightarrow\sum_i\alpha_i=0,\\
\partial_{e_i}\mathcal L=0&\Rightarrow e_i=\alpha_i/\gamma, &
\partial_{\alpha_i}\mathcal L=0&\Rightarrow y_i=w^T\phi_i+b+e_i.
\end{align*}
代入最后一式：$y_i=\sum_jK_{ij}\alpha_j+b+\alpha_i/\gamma$。再加上乘子和为零，得到
\[
\boxed{\begin{bmatrix}K+\gamma^{-1}I&\mathbf1\\\mathbf1^T&0\end{bmatrix}
\begin{bmatrix}\alpha\\b\end{bmatrix}=\begin{bmatrix}y\\0\end{bmatrix}},
\qquad f(x)=\sum_i\alpha_iK(x_i,x)+b.
\]
\ans{LS-SVM回归最小化范数与平方残差之和，等式为$y=f+e$。消去$w,e$后求上述$(n+1)$阶系统；预测是连续值，等式乘子可正可负。}
\qcheck{上方$n$行给$n$个训练等式，最后一行来自bias导数；维度和未知数都为$n+1$。回代检验$e=y-K\alpha-b\mathbf1=\alpha/\gamma$以及$\mathbf1^T\alpha=0$。}
\pitfall{不能在预测中额外乘分类标签$y_i$；不能给等式乘子加$\alpha_i\ge0$；不能漏掉bias导数。}
\transfer{换核或数据只改变$K,y$；换$\gamma$改变对角$1/\gamma$。若题目明确不允许bias，则删掉$b$及对应末行末列，不能继续套同一系统。}

\problem{Q47}{两样本LS-SVM：增广消元与完整回代}{背面第4栏 L01、L02；背面第5栏 L03；正面第1栏 M03}
\qtext{训练样本为$(x_1,y_1)=(0,0)$、$(x_2,y_2)=(1,2)$。取线性核$K(x,z)=xz$、$\gamma=1$，包含不惩罚的偏置$b$。求$K$、增广系统、$\alpha,b,w$、$f(2)$、训练残差及最优目标值。必须展示行变换，不能仅写矩阵求逆。}
\route{圈“线性核、$\gamma$、平方回归、包含bias”。L02给系统，L03是同题的压缩版；M03指导行变换。目标连续，不能套分类的$y^T\alpha=0$。}
\given{$n=2$，$y=(0,2)^T$，未知$q=(\alpha_1,\alpha_2,b)^T$，$\phi(x)=x$，$w$为标量。}
\solve
先逐项算核：$K_{11}=0,K_{12}=K_{21}=0,K_{22}=1$。对角加1，补偏置行列：
\[
K=\begin{bmatrix}0&0\\0&1\end{bmatrix},\quad
A=\begin{bmatrix}1&0&1\\0&2&1\\1&1&0\end{bmatrix},\quad Aq=\begin{bmatrix}0\\2\\0\end{bmatrix}.
\]
增广消元按以下顺序进行，右端必须跟随整行变化：
\[
\left[\begin{array}{ccc|c}1&0&1&0\\0&2&1&2\\1&1&0&0\end{array}\right]
\xrightarrow{R_3\leftarrow R_3-R_1}
\left[\begin{array}{ccc|c}1&0&1&0\\0&2&1&2\\0&1&-1&0\end{array}\right]
\xrightarrow{R_2\leftrightarrow R_3}
\left[\begin{array}{ccc|c}1&0&1&0\\0&1&-1&0\\0&2&1&2\end{array}\right].
\]
\[
\xrightarrow{R_3\leftarrow R_3-2R_2}
\left[\begin{array}{ccc|c}1&0&1&0\\0&1&-1&0\\0&0&3&2\end{array}\right].
\]
从末行回代：$3b=2$，所以$b=2/3$；第二行$\alpha_2-b=0$给$\alpha_2=2/3$；第一行$\alpha_1+b=0$给$\alpha_1=-2/3$。
\[
w=\sum_i\alpha_ix_i=2/3,\qquad f(x)=\frac23x+\frac23,\qquad f(2)=2.
\]
训练预测为$(2/3,4/3)^T$，故$e=y-f=(-2/3,2/3)^T$。目标值
\[
J=\tfrac12(2/3)^2+\tfrac12\big[(-2/3)^2+(2/3)^2\big]=\frac23.
\]
\ans{$\alpha=(-2/3,2/3)^T,b=w=2/3,f(2)=2$，训练残差$(-2/3,2/3)^T$，最优目标值$2/3$。}
\qcheck{$Aq=(0,2,0)^T$；$\alpha_1+\alpha_2=0$；$\gamma e=\alpha$；对$w$求导条件$w-\sum\alpha_i x_i=0$。四组条件全满足，凸问题中足以确认最优。}
\pitfall{最后一行右端为0；不能给$b$的右下角写1。$f(2)$是新点预测，不能用第二个训练残差代替。负乘子在这里正常。}
\transfer{数据顺序可交换，但$K$的行列、$y$、$\alpha$顺序必须同步；核换成非线性后仍按同一框架求解。}

\problem{Q48}{改变正则权重：哪些量要重算}{背面第5栏 L03；背面第4栏 L02；正面第4栏 B03}
\qtext{沿用Q47的两个样本和线性核，先把$\gamma$改为2，再对一般$\gamma>0$解出$w,b,\alpha$。比较训练平方误差和权重范数随$\gamma$的变化，并判断“$\gamma$越大测试必越准”是否正确。}
\route{圈“同一数据、改变$\gamma$、训练与测试”。L03提示只改$K+\gamma^{-1}I$的对角，但必须重新求整个系统；B03区分拟合和泛化。}
\given{$K=\operatorname{diag}(0,1)$，$y=(0,2)^T$；$\gamma$是平方残差惩罚的权重，非核宽度。令$\alpha_2=a$。}
\solve
由末行$\alpha_1+\alpha_2=0$得$\alpha_1=-a$。前两行分别是
\[
-a/\gamma+b=0,\qquad (1+1/\gamma)a+b=2.
\]
第一式给$b=a/\gamma$，代第二式得到$a(1+2/\gamma)=2$，所以
\[
\alpha=\left(-\frac{2\gamma}{\gamma+2},\frac{2\gamma}{\gamma+2}\right)^T,
\quad b=\frac2{\gamma+2},\quad w=\frac{2\gamma}{\gamma+2}.
\]
当$\gamma=2$：$\alpha=(-1,1)^T,b=1/2,w=1$；训练预测$(1/2,3/2)$，残差$(-1/2,1/2)$，$f(2)=5/2$。
\[
\sum e_i^2=\frac8{(\gamma+2)^2},\qquad
\|w\|^2=\frac{4\gamma^2}{(\gamma+2)^2}.
\]
本数据中前者随$\gamma$增大而减小，后者增大。$\gamma\to0^+$时$w\to0,b\to1$，接近常数均值预测；$\gamma\to\infty$时$w\to2,b\to0$，趋向插值。这些极限来自本题可插值数据，不是任意数据都能零残差。
\ans{$\gamma=2$时$w=1,b=1/2,\alpha=(-1,1)^T$。更大$\gamma$更重视训练平方误差，可能使用更大范数；测试表现必须独立验证。}
\qcheck{代回新矩阵的首行$\alpha_1/2+b=-1/2+1/2=0$；次行$(3/2)\alpha_2+b=2$；$e=\alpha/2$。}
\pitfall{不要把$1/\gamma$写成$\gamma$。一般问题可证明最优总训练SSE随该权重增加不增，但不能据此断言每一个样本的残差都单调缩小。}
\transfer{换数据后仍保留“算核、加对角、加bias、消元、回代”的流程；具体单调公式和极限必须重新推导。核宽度变化还会改变非对角元素，不能只修改对角。}

\problem{Q49}{找出LS-SVM推导中的六类错误}{背面第4栏 L01、L02；背面第5栏 L03；背面第1栏 S02}
\qtext{同学拟合连续目标$y_i\in\R$，声称LS-SVM必须穿过全部样本。他写
\[
\mathcal L=\tfrac12\|w\|^2+\tfrac\gamma2\sum e_i^2-\sum\alpha_i(w^T\phi_i+b+e_i-y_i),
\]
却得到$w=\sum\alpha_i y_i\phi_i$、$e_i=\gamma\alpha_i$、$0\le\alpha_i\le C$，并列出
\[
\begin{bmatrix}K-\gamma I&\mathbf1\\y^T&1\end{bmatrix}
\begin{bmatrix}\alpha\\b\end{bmatrix}=\begin{bmatrix}y\\0\end{bmatrix}.
\]
逐类改正，并解释为什么把未知向量改成$[b;\alpha]$而不改矩阵也不对。}
\route{圈“连续目标、等式、LS-SVM”，回L01按给出的L重新求导；不要从错误结果倒猜。S02只用作与分类公式的对照。}
\given{$K_{ij}=\phi_i^T\phi_j$，$\gamma>0$；题中L本身正确，错误从导数及解释开始。$C$并非本模型给定参数。}
\solve
\begin{enumerate}
\item 等式$y_i=f_i+e_i$允许非零残差，不能推成必须插值。
\item $\partial_w\mathcal L=w-\sum\alpha_i\phi_i$，故$w=\sum\alpha_i\phi_i$，没有分类标签因子。
\item $\partial_{e_i}\mathcal L=\gamma e_i-\alpha_i$，故$e_i=\alpha_i/\gamma$；不是$\gamma\alpha_i$。
\item 等式乘子$\alpha_i\in\R$，没有非负或盒约束；不能引入不存在的$C$。
\item 代入$e$产生$+\gamma^{-1}I$，既不是减号，也不是$\gamma I$。
\item bias导数给$\sum_i\alpha_i=0$，所以末行是$[\mathbf1^T,0]$，不是$[y^T,1]$。
\end{enumerate}
\[
\text{正确系统：}\quad
\begin{bmatrix}K+\gamma^{-1}I&\mathbf1\\\mathbf1^T&0\end{bmatrix}
\begin{bmatrix}\alpha\\b\end{bmatrix}=\begin{bmatrix}y\\0\end{bmatrix}.
\]
未知向量换成$[b;\alpha]$，若也把bias方程排第一行，则应同时得到
\[
\begin{bmatrix}0&\mathbf1^T\\\mathbf1&K+\gamma^{-1}I\end{bmatrix}
\begin{bmatrix}b\\\alpha\end{bmatrix}=\begin{bmatrix}0\\y\end{bmatrix}.
\]
\ans{等式允许残差；$w=\sum\alpha_i\phi_i,e=\alpha/\gamma,\alpha\in\R^n$；左上$K+I/\gamma$，末行$[\mathbf1^T,0]$。未知顺序变动必须匹配矩阵列，方程顺序变动必须匹配行及右端。}
\qcheck{Q47的$\alpha=(-2/3,2/3)$直接否定“全非负”。正确系统维度为$(n+1)\times(n+1)$，每行能还原对应导数或训练等式。}
\pitfall{“核方法”并不意味着所有模型的预测都含$y_i$；“SVM”后缀也不意味着所有乘子约束一致。}
\transfer{若选择相反的Lagrangian乘子符号，可整体重新定义乘子并同步改所有式子；不能只反转某一条导数或预测符号。}

\problem{Q50}{三样本非线性核：矩阵秩低仍可唯一求解}{背面第2栏 K01；背面第4栏 L02；背面第5栏 L03；正面第1栏 M03}
\qtext{样本$x=(-1,0,1)$、$y=(1,0,1)$，核$K(x,z)=(xz)^2$，$\gamma=1$，包含bias。说明核有效，计算Gram矩阵、完整系统及解，写出$f(x)$、$f(2)$和训练残差。不能因为$K$奇异就停止求解。}
\route{圈“连续目标、非线性核、平方残差”；用K01确认特征映射，用L02搭系统。不是FLN的无正则精确插值；本题带范数惩罚。}
\given{可取$\phi(x)=x^2\in\R$，所以特征权重$w$为标量。未知按$(\alpha_1,\alpha_2,\alpha_3,b)$排列。}
\solve
$K_{ij}=x_i^2x_j^2$，故$K=vv^T$，$v=(1,0,1)^T$；任意$c$有$c^TKc=(v^Tc)^2\ge0$，核有效。
\[
K=\begin{bmatrix}1&0&1\\0&0&0\\1&0&1\end{bmatrix},\quad
[A|r]=\left[\begin{array}{cccc|c}2&0&1&1&1\\0&1&0&1&0\\1&0&2&1&1\\1&1&1&0&0\end{array}\right].
\]
依次做$R_1\leftarrow R_1-R_3$，再$R_3\leftarrow R_3-R_1$、$R_4\leftarrow R_4-R_1$（此时用新$R_1$），再$R_4\leftarrow R_4-R_2$：
\[
\left[\begin{array}{cccc|c}1&0&-1&0&0\\0&1&0&1&0\\0&0&3&1&1\\0&0&2&-1&0\end{array}\right]
\xrightarrow{R_3\leftrightarrow R_4,\ R_4\leftarrow R_4-\frac32R_3}
\left[\begin{array}{cccc|c}1&0&-1&0&0\\0&1&0&1&0\\0&0&2&-1&0\\0&0&0&5/2&1\end{array}\right].
\]
回代得$b=2/5$，$\alpha_3=b/2=1/5$，$\alpha_2=-b=-2/5$，$\alpha_1=\alpha_3=1/5$。
\[
w=\sum\alpha_ix_i^2=2/5,\quad
f(x)=\tfrac25x^2+\tfrac25,\quad f(2)=2.
\]
训练预测为$(4/5,2/5,4/5)$，残差$(1/5,-2/5,1/5)$。目标$J=\frac12(2/5)^2+\frac12(1+4+1)/25=1/5$。
\ans{$\alpha=(1/5,-2/5,1/5)^T,b=2/5$；$f(x)=2x^2/5+2/5$，$f(2)=2$。$K$秩1，但$K+I$正定，含bias系统仍唯一可解。}
\qcheck{$\sum\alpha_i=0$；残差等于$\alpha$；$K\alpha+b\mathbf1+\alpha=y$。核预测与显式一维特征预测得到相同$f$。}
\pitfall{不能把$K(x,z)$写成$(x-z)^2$，后者不是本题核；不能去掉中间$x=0$样本，它仍影响bias及残差。}
\transfer{换成另一个有效核时，不一定能手写有限维$\phi$，但Gram系统流程不变。增加样本会同时增加未知乘子和训练方程。}

\problem{Q51}{唯一解、二次优化与“最大”辨析}{背面第5栏 L03、D06；正面第1栏 M02；背面第1栏 S01}
\qtext{设$n\ge1$，$K$为对称半正定Gram矩阵，$\gamma>0$。证明$H=K+\gamma^{-1}I$正定，进而说明LS-SVM含bias系统可逆。判断并改正：“分块矩阵右下为0，所以不可逆”；“用线性方程组训练，所以目标是线性的”；“老师提到最大，因此所有模型都最大化损失”；“LS-SVM所有乘子都非零，所以永远稠密”。}
\route{圈“PSD、正定、唯一、目标方向、通常”。L03给条件，M02帮助区分驻点与极值；S01解释最大间隔的具体含义。此题是条件与证明，不需要虚构数值。}
\given{$H\in\R^{n\times n}$；令齐次系统未知为$(a,c)$，$a\in\R^n,c\in\R$。用齐次系统只有零解判断方阵可逆。}
\solve
任意$a\ne0$，有
\[
a^THa=a^TKa+\gamma^{-1}\|a\|^2\ge\gamma^{-1}\|a\|^2>0.
\]
因此$H$正定。考虑$Ha+c\mathbf1=0$、$\mathbf1^Ta=0$；第一式左乘$a^T$得
\[
a^THa+c\,a^T\mathbf1=0\Rightarrow a^THa=0\Rightarrow a=0.
\]
代回得$c\mathbf1=0$，由于$n\ge1$，所以$c=0$。于是含bias方阵只有零齐次解，故可逆，任意右端有唯一解。右下一个元素为0不代表整矩阵奇异；该分块矩阵也不必正定。

LS-SVM原目标仍为范数平方和残差平方的\textbf{二次凸优化}；等式约束使最优条件可整理成线性系统。hard SVM最大化几何间隔等价于最小化范数；LMS、BP平方损失、LS-SVM均采用相应最小化问题。用户提供的“最高／最大”未说明具体量，不能当作所有题的方向。

最后，$\alpha_i=\gamma e_i$，一般残差不全为0，因此标准LS-SVM通常有很多非零系数；特殊数据可出现零系数，不能把“通常稠密”写成“每个系数必非零”。
\ans{$K\succeq0,\gamma>0,n\ge1$保证$H\succ0$及含bias系统可逆。线性驻点系统源于二次优化；最大或最小必须看题目目标；LS-SVM通常稠密但允许零系数。}
\qcheck{Q47的$K$奇异却有唯一解，符合结论。若所有$y_i$相同，可取$w=0,b=y_i,e=0,\alpha=0$，直接反驳“所有乘子都非零”。}
\pitfall{不能从$K$对称就推出PSD；也不能省略$\gamma>0$后继续使用正定证明。不要把所有无约束驻点都叫全局最小。}
\transfer{若题给无效核，需先查假设；若不含bias，系统变为$H\alpha=y$，正定性直接保证唯一。模型比较仍按任务、目标、约束、求解及限制五个维度组织。}

\clearpage\section*{混合自测：先遮住后面的解析}\addcontentsline{toc}{section}{混合自测题目（Q52--Q56）}
这五题均为自拟复习题，不表示考试比例或时限。先只看cheatsheet独立完成，圈出关键词并写下查表编号；从下一组Q52解析开始再核对。每题同时考查定位、计算与解释。

\subsection*{Q52 同一个输入，三种更新为什么不同}
查表：正面第2栏P01；第3栏A01、B01；第4栏B02。

单样本$x=1$、目标$z=1$，初始$w=.2,b=.1,\eta=.1$。分别把该单元解释为：(A) $\hat y=H(wx+b),H(0)=1$的感知机；(B) $a=wx+b$的ADALINE，$E=\frac12(z-a)^2$；(C) $a=\sigma(wx+b)$的单个输出神经元，仍用半平方误差。求三者一次同步更新后的$w,b$，说明为何仅凭“更新权重”不能选公式。

\subsection*{Q53 有核系数就一定是最优分类器吗}
查表：背面第1栏S01、S02；第2栏K01。

线性核，训练$x=(-2,2)$、$y=(-1,1)$。有人交出$\alpha_1=\alpha_2=1/8,b=1$，并称两点都是自由支持向量。检查可行性，纠正$b$，再求$w$、总间隔和primal/dual值。若将$(w,b)$再乘2，哪些量保持不变？这里按hard-margin讨论。

\subsection*{Q54 加对角能否把无效核变成原来的有效核}
查表：背面第2栏K01；第4栏L02；第5栏L03。

拟用$K=\left[\begin{smallmatrix}1&2\\2&1\end{smallmatrix}\right]$作为两点Gram矩阵，$y=(0,2)^T,\gamma=1/2$。先找一个向量证明该Gram无效，再检查$H=K+2I$是否正定。若形式上求带bias系统，会得到什么解？能否据此宣布原$K$是有效核？

\subsection*{Q55 改变管宽后，原来的KKT证据还成立吗}
查表：背面第3栏V01、V02。

一条样本$x=0,y=3$，线性核，$C=1$。先取$\varepsilon=0$，候选$w=0,b=3,\xi=\xi^*=0,\alpha=\alpha^*=2/5$。逐项检查KKT。再改$\varepsilon=1/2$而保持这些变量不变：原预测是否仍最优？原乘子是否仍是有效KKT证据？给出可行的修正乘子。

\subsection*{Q56 训练目标最优，能否直接宣布泛化最好}
查表：背面第4栏L01；第5栏L03；正面第4栏B03。

训练样本$(0,0),(1,2)$，线性LS-SVM，$\gamma=1$，bias不惩罚。候选A为$f_A(x)=2x$，候选B为$f_B(x)=2x/3+2/3$。计算两者训练目标。另有一个此前未用来拟合的测试样本$(2,4)$，比较该点平方预测误差。解释两种排序为什么可相反，以及这个单点能支持多强的泛化结论。

\problem{Q52}{混合自测解析：同分数、不同学习信号}{正面第2栏P01；第3栏A01、B01；第4栏B02}
\qtext{数据$x=z=1,w=.2,b=.1,\eta=.1$。对硬阈值感知机、线性ADALINE、sigmoid输出加半平方误差三种设定，各作一步同步更新。初始净输入均为$u=.3$。}
\route{“硬类别”指P01；“线性连续输出的平方残差”指A01；“sigmoid加链式导数”指B02。三者有同一个$u$，但预测和损失对参数的依赖不同。}
\given{$w,b$均为标量；$H(u)=1$当$u\ge0$；$\sigma(u)=1/(1+e^{-u})$；每种情况独立从相同旧参数开始，不能顺接更新。}
\solve
(A) 感知机：$\hat y=H(.3)=1$，$z-\hat y=0$，所以$w^+=.2,b^+=.1$。

(B) ADALINE：$a=.3$，线性残差$e=.7$，所以$\Delta w=.1(.7)(1)=.07$、$\Delta b=.07$，得到$w^+=.27,b^+=.17$。

(C) sigmoid输出：$a=\sigma(.3)\approx.57444252$。使用课件负导数约定，
\[
\delta=(z-a)a(1-a)\approx.10403106,
\quad\Delta w=\eta\delta x\approx.01040311,\quad\Delta b\approx.01040311.
\]
故$w^+\approx.21040311,b^+\approx.11040311$。这里是一个输出单元，无需额外编造隐藏层。
\ans{三者新参数分别为$(.2,.1)$、$(.27,.17)$、约$(.21040311,.11040311)$。相同输入和学习率不能确定更新式；必须核对输出激活及训练误差信号。}
\qcheck{感知机分类已经正确，零更新合理；LMS更新后$a=.44$，半平方误差由$.245$降到$.1568$；sigmoid情形通过直接前向损失的有限差分可验证梯度符号。}
\pitfall{C不能直接套$e=.7$，也不能只用$z-a$而漏$a(1-a)$。A不能因$u=.3\ne1$就误以为分类错误。}
\transfer{改为线性输出时删激活导数；目标若改成别的损失须重新链式求导。若新增隐藏层，还要用同一旧输出权重反传。}

\problem{Q53}{混合自测解析：偏置、间隔与最优性一起检查}{背面第1栏S01、S02；背面第2栏K01}
\qtext{线性核，$x=(-2,2),y=(-1,1),\alpha=(1/8,1/8)$，候选$b=1$。检查并纠正hard-margin解，再分析正比例缩放。}
\route{圈“hard-margin、给乘子、检查偏置”。S02从系数恢复$w$并检查KKT，S01检查几何间隔；只看到乘子非负不足以证明最优。}
\given{$w=\sum\alpha_i y_ix_i$为标量；函数间隔$m_i=y_i(wx_i+b)$；无软间隔$C$，题中的“自由SV”应理解为hard中$\alpha_i>0$。}
\solve
先算$w=(1/8)(-1)(-2)+(1/8)(1)(2)=1/2$。若$b=1$，则$g=(0,2)$，函数间隔$(0,2)$，第一点违反$m\ge1$，因此不可行。

两点乘子均正，应满足$m=1$。用任一点恢复：
\[
b=y_1-wx_1=-1-(1/2)(-2)=0,
\quad y_2-wx_2=1-1=0.
\]
正确解为$w=1/2,b=0$，两间隔均1，总宽$2/|w|=4$。primal为$\frac12w^2=1/8$；dual为
\[
\sum\alpha_i-\frac12\left(\sum\alpha_i y_i x_i\right)^2
=\frac14-\frac18=\frac18.
\]
非负、$\sum\alpha_i y_i=0$、原可行、驻点、互补均满足，且原对偶同值。

把$(w,b)$乘2，变为$(1,0)$：决策边界仍$x=0$，分类不变，几何间隔仍2；函数间隔从1变2。此时数据不再处在规范的$g=\pm1$支撑面上，不能把$2/\|w\|=2$当本数据的最大总宽。
\ans{错误是$b=1$；正确$b=0,w=1/2$，总间隔4，原/对偶最优值均$1/8$。正比例缩放保留边界和几何距离，改变函数间隔与规范化。}
\qcheck{直接点到面距离$|x_i|=2$，两侧总宽4，与规范公式一致。}
\pitfall{“已给$\alpha$”不等于给定答案正确；必须连同$b$检查全部约束。}
\transfer{soft-margin求$b$要改用$0<\alpha<C$；核化后通常不显式求$w$，用$\sum\alpha_i y_iK_{is}$恢复$b$。}

\problem{Q54}{混合自测解析：能解方程不等于核合法}{背面第2栏K01；第4栏L02；第5栏L03}
\qtext{$K=\left[\begin{smallmatrix}1&2\\2&1\end{smallmatrix}\right]$、$\gamma=1/2,y=(0,2)^T$。证明$K$无效，检查$H=K+2I$并形式求解，然后指出不能得出的结论。}
\route{圈“Gram有效性、加对角、能否宣布”。先K01检查原核，再L02做代数；L03的唯一性充分条件不能反向推回原核必有效。}
\given{$c=(1,-1)^T$用作PSD反例；$H=\left[\begin{smallmatrix}3&2\\2&3\end{smallmatrix}\right]$，未知仍$[\alpha_1,\alpha_2,b]$。这里只作错误诊断，不将无效核纳为标准LS-SVM。}
\solve
$c^TKc=-2<0$，所以$K$不是PSD，不能是实内积特征的Gram矩阵。$H$的特征值是$3+2=5$和$3-2=1$，均正，故本矩阵$H$确实正定。

形式系统为
\[
\left[\begin{array}{ccc|c}3&2&1&0\\2&3&1&2\\1&1&0&0\end{array}\right].
\]
末行给$\alpha_1=-\alpha_2$；令$\alpha_2=a$，前两行化为$-a+b=0$、$a+b=2$，因此$a=1,b=1$，即$\alpha=(-1,1)^T$。

这说明一个给定代数系统可解；并没有改变原$K$的负二次型。若把$H$本身重新定义为另一个有限集合上的核矩阵，那已是换模型，也仍要说明完整核函数及适用域，不能声称原$K$被“证明有效”。
\ans{原$K$无效；$H$正定且形式系统解为$\alpha=(-1,1),b=1$。可逆/可解性质不能替代有效核要求。}
\qcheck{$K\alpha+b\mathbf1=(2,0)^T$，形式残差$(-2,2)^T=\alpha/\gamma$；代数条件成立，但$c^TKc=-2$仍未消失。}
\pitfall{不能从一个有限点集的修补矩阵PSD推出原核函数对所有输入有效；更不能从对称性单独推出PSD。}
\transfer{若换成有效$K\succeq0$且$\gamma>0$，L03的保证可直接使用。若Gram计算漏了特征映射的$\sqrt2$项，先重算正确核，不应任意改数字掩盖错误。}

\problem{Q55}{混合自测解析：预测最优与乘子证据不同}{背面第3栏V01、V02}
\qtext{一条样本$x=0,y=3$，线性SVR，$C=1$；候选$w=0,b=3,\xi=\xi^*=0,\alpha=\alpha^*=2/5$。先验$\varepsilon=0$，再换为$1/2$；判断原预测和原乘子是否均可保留。}
\route{圈“$\varepsilon=0$、两乘子同正、换管宽”。V02明示零管宽边界需重判；要区分原始变量的最优性和给定乘子的KKT有效性。}
\given{$r=y-f=0$，$d=\alpha-\alpha^*=0$，核$K(0,0)=0$；由驻点可取松弛非负约束的乘子$\eta=\eta^*=C-2/5=3/5$。}
\solve
当$\varepsilon=0$：两个误差约束都是$0\le0$，松弛非负；乘子在盒内；$\sum d=0$，$w=dx=0$。两条误差约束互补为$(2/5)\cdot0=0$，松弛互补同为0。原目标为0，对偶目标也为0，故KKT与最优性均成立。

改$\varepsilon=1/2$后，原始预测仍有$r=0$、零松弛、$w=0$，目标0已经是非负目标的下界，因此\textbf{预测仍最优}。然而误差约束有严格余量$1/2$，原乘子的互补乘积变为$(2/5)(1/2)=1/5\ne0$，所以\textbf{原乘子失效}。改为$\alpha=\alpha^*=0,\eta=\eta^*=1$即可满足全部条件。
\ans{$\varepsilon=0$允许同一点两乘子同时正；增到$1/2$时原预测仍最优，但必须改乘子，取两乘子0是一个正确选择。}
\qcheck{增管宽而保留原乘子会给dual值$-\varepsilon(\alpha+\alpha^*)=-2/5$，与primal值0不相等；改成零乘子后二者同为0。}
\pitfall{不能因为预测没有改变就跳过互补检查；也不能把“$\varepsilon>0$不能两乘子同正”无条件套在零管宽。}
\transfer{若残差不为0，增管宽还可能改变原始最优解；重新算管外损失、约束和KKT。无自由SV时使用可行bias区间，不随意套边界等号。}

\problem{Q56}{混合自测解析：目标与泛化评价可以排序相反}{背面第4栏L01；第5栏L03；正面第4栏B03}
\qtext{训练$(0,0),(1,2)$、线性LS-SVM、$\gamma=1$。比较A:$f=2x$与B:$f=2x/3+2/3$的训练目标；再比较未参与拟合的测试点$(2,4)$的平方误差，并解释结论边界。}
\route{圈“范数加平方误差、测试样本”。L01决定训练目标，B03决定如何解释测试结果。单看训练残差为0不能宣布A的LS目标更优。}
\given{A有$w_A=2,b_A=0$；B有$w_B=b_B=2/3$。训练目标$J=\frac12w^2+\frac12\sum e_i^2$；测试指标为$(y-f)^2$，题目未在此加$1/2$。}
\solve
A在两个训练点完全插值，$e=(0,0)$，但范数代价$\frac12(2)^2=2$，所以$J_A=2$。

B的训练残差$(-2/3,2/3)$，所以
\[
J_B=\frac12\frac49+\frac12\left(\frac49+\frac49\right)=\frac23<2.
\]
B是Q47已由最优条件验证的该训练问题解。对新点$x=2$，$f_A(2)=4$，平方误差0；$f_B(2)=2$，平方误差$(4-2)^2=4$。于是训练正则目标偏好B，但给定测试点偏好A；没有矛盾，因为评价对象和依据不同。

一个测试点只能报告这个点的结果，不能证明A对总体分布一定更好。模型和超参数应使用训练/验证过程选择；测试数据保留最终评价，不能在观察结果后反复调参还把它当独立测试。
\ans{$J_A=2,J_B=2/3$；测试平方误差分别0和4。训练目标最优不保证每个未见点更准，单点测试也不足以确定整体泛化优劣。}
\qcheck{B的残差与$\alpha/\gamma$一致；A虽然训练误差零，但$w=2$且若$e=0$则$\alpha=0$，不满足$w=\sum\alpha_i x_i$，所以不是该正则问题最优解。}
\pitfall{别把“平方误差”默认都带$1/2$；本题明确区分训练目标与测试指标。也不要把用户提供的“最大”误读成要最大化这些损失。}
\transfer{换$\gamma$会改变训练权衡；换测试分布会改变泛化排序。保持“先定义指标、再计算、再限定结论”的步骤，不能仅凭模型名称比较好坏。}

\end{document}
